\subsection{正线性泛函}

定义 2. --- 设 A 是对合代数。A 上的线性泛函 \(\lambda\) 称为正的，如果 \(\lambda(a^{*}a) \in \mathbf{R}_{+}\) 对所有 \(a \in \mathrm{A}\) 都成立。

若 A 是对合巴拿赫代数，我们以 \(A_{+}^{\prime}\) 表示 A 上连续正线性泛函的集合。

设 A 为对合巴拿赫代数。集合 \(A_{+}^{\prime}\) 是 A 上 C-线性泛函所在实向量空间中的尖凸锥。

引理 2. --- 设 A 为对合巴拿赫代数，且 \(\lambda\) 是 A 上的正线性泛函。

\begin{enumerate}
\def\labelenumi{\alph{enumi})}
\tightlist
\item
  对 A 中任意 \(a\) 和 \(b\)，有 \(\lambda(a^{*}b)=\overline{\lambda(b^{*}a)}\)，且
\end{enumerate}

\[
\left| \lambda \left(a ^ {*} b\right) \right| ^ {2} \leqslant \lambda \left(a ^ {*} a\right) \lambda \left(b ^ {*} b\right);
\]

\begin{enumerate}
\def\labelenumi{\alph{enumi})}
\setcounter{enumi}{1}
\tightlist
\item
  若 A 含单位元，则线性泛函 \(\lambda\) 连续，且其范数等于 \(\lambda(1)\)。
\end{enumerate}

映射 \((a,b)\mapsto\lambda(a^{*}b)\) 是 A 上的正埃尔米特形式；因此，对 A 中任意 a 和 b，它满足 \(\lambda(a^{*}b)=\overline{\lambda(b^{*}a)}\) 以及 \(|\lambda(a^{*}b)|^{2}\leqslant\lambda(a^{*}a)\lambda(b^{*}b)\)（EVT, V, p. 2, 注以及 p. 3, 命题 2）。

证明 \(b\)。设 \(a \in A\) 是范数 \(< 1\) 的埃尔米特元。\(a\) 的谱半径小于 \(\|a\|\)，所以 \(a\) 的谱包含于 C 的开单位圆盘 D 中（I, p. 24, 定理 1）。设 \(f\) 是 D 上满足 \(f(z)^{2} = 1 - z\) 且对所有 \(z \in D\) 成立的全纯函数（V, p. 435, 引理 1）。对元 \(a\) 和函数 \(f\) 应用全纯函数演算。元 \(b = f(a)\) 满足 \(b^{2} = 1 - a\)（I, p. 74, 定理 5）。由 V, p. 435, 命题 3，还有 \(f(a)^{*} = f^{*}(a^{*}) = f(a)\)，故 \(b\) 是埃尔米特元。于是 \(\lambda(1 - a) = \lambda(b^{*}b) \geqslant 0\)，从而 \(\lambda(a) \leqslant \lambda(1)\)。

现在设 \(b \in A\) 且范数 \(< 1\)。元 \(b^*b\) 是埃尔米特元且 \(\|b^*b\| < 1\)，因此在应用 \(a\) )（取 \(a = 1\)）的情形，得到

\[
| \lambda (b) | ^ {2} \leqslant \lambda (1) \lambda (b ^ {*} b) \leqslant \lambda (1) ^ {2},
\]

当 \(b = 1\) 时取等号。因此线性泛函 \(\lambda\) 连续，且其范数等于 \(\lambda(1)\)。断言 \(b\) ) 得证。

例。--- 设 X 是紧拓扑空间，且 A 是星代数 \(\mathcal{C}(\mathrm{X})\)。A 上的正线性泛函与 X 上的正测度相对应（INT, III, p. 52, § 1, n° 6, 定理 1）。

引理 3. --- 设 A 是具有逼近单位元的含单位元对合巴拿赫代数（I, p. 120, 定义 7）。设 λ 是 A 上的连续正线性泛函。

\begin{enumerate}
\def\labelenumi{\alph{enumi})}
\item
  对 A 中任意 a，有 \(\lambda(a^{*})=\overline{\lambda(a)}\) 以及 \(|\lambda(a)|^{2}\leqslant\|\lambda\|\lambda(a^{*}a)\)；
\item
  设 \(\widetilde{\mathrm{A}}\) 是通过向 A 添加单位元得到的对合代数，e 是其单位元。存在 \(\widetilde{\lambda}\) 在 \(\widetilde{\mathrm{A}}\) 上的连续正线性泛函，它延拓 \(\lambda\)，且满足 \(\widetilde{\lambda}(e) = \|\lambda\|\)；
\item
  对 A 中任意 \(a\) 和 \(b\)，有 \(|\lambda(b^{*}ab)|\leqslant\|a\|\lambda(b^{*}b)\)。
\end{enumerate}

证明 \(a\) )。设 \(\mathfrak{F}\) 是 A 的逼近单位元，且 \(a \in \mathrm{A}\)。利用引理 2、\(a\) ) 以及逼近单位元的定义，可得

\[
\lambda (a ^ {*}) = \lim _ {f, \mathfrak {F}} \lambda (f a ^ {*}) = \lim _ {f, \mathfrak {F}} \overline {{\lambda (a f ^ {*})}} = \lim _ {f, \mathfrak {F}} \overline {{\lambda ((f a ^ {*}) ^ {*})}} = \overline {{\lambda (a)}},
\]

因此（同上）

\[
| \lambda (a) | ^ {2} = \lim _ {f, \mathfrak {F}} | \lambda (f a) | ^ {2} \leqslant \lambda (a ^ {*} a) \operatorname * {limsup} _ {f, \mathfrak {F}} \lambda (f f ^ {*}) \leqslant \| \lambda \| \lambda (a ^ {*} a),
\]

因为 \(\mathfrak{F}\) 是 A 的单位球上的滤子。

证明 \(b\) )。不妨设 \(\lambda\) 非零，进而设 \(\| \lambda \| = 1\)。对 \(a \in \mathrm{A}\) 和 \(z \in \mathbf{C}\)，定义 \(\widetilde{\lambda}(a + z \cdot e) = \lambda(a) + z\)。映射 \(\widetilde{\lambda}\) 是 \(\widetilde{\mathrm{A}}\) 上延拓 \(\lambda\) 的连续线性泛函，并满足 \(\widetilde{\lambda}(e) = 1\)。它是正的：对任意 \(a \in \mathrm{A}\) 和 \(z \in \mathbf{C}\)，根据 \(a\) )，计算得

\[
\begin{array}{r l} \widetilde {\lambda} ((a + z \cdot e) ^ {*} (a + z \cdot e)) & = \lambda (a ^ {*} a) + \overline {{z}} \lambda (a) + z \lambda (a ^ {*}) + | z | ^ {2} \\ & = | z + \lambda (a) | ^ {2} + \lambda (a ^ {*} a) - | \lambda (a) | ^ {2} \geqslant 0. \end{array}
\]

最后，设 \(b\) 是 A 的一个元。映射 \(a \mapsto \widetilde{\lambda}(b^{*}ab)\) 是含单位元对合巴拿赫代数 \(\widetilde{A}\) 上的正线性泛函。因此它连续，且范数等于其在 \(e\) 处的值（引理 2，\(b\) )），而该值等于 \(\lambda(b^{*}b)\)，从而得到断言 \(c\) )。

命题 4. — 设 A 为对合巴拿赫代数。

\begin{enumerate}
\def\labelenumi{\alph{enumi})}
\item
  对每个希尔伯特空间 E、每个对合代数态射 \(\varphi\colon A\to\mathcal{L}(E)\) 以及每个向量 \(x\in E\)，在 A 上由 \(\lambda(a)=\langle x|\varphi(a)x\rangle\) 定义的线性泛函是连续正线性泛函；
\item
  设 \(\lambda \in \mathrm{A}_{+}^{\prime}\)。映射 \(f\) 从 \(\mathrm{A} \times \mathrm{A}\) 到 \(\mathbf{C}\)，由 \(f(a, b) = \lambda (a^{*}b)\) 对所有 \((a, b) \in \mathrm{A} \times \mathrm{A}\) 定义，是 \(\mathrm{A}\) 上的普遍正核。
\end{enumerate}

证明 \(a\) )。对每个 \(a \in \mathrm{A}\)，有

\[
\lambda (a ^ {*} a) = \langle x \mid \varphi (a ^ {*} a) x \rangle = \| \varphi (a) x \| ^ {2} \geqslant 0
\]

所以 \(\lambda\) 是 A 上的正线性泛函；由于 \(\varphi\) 连续，它也是连续的（I, p. 104, 命题 2）。

证明 \(b\) )。函数 \(f\) 连续。对每个 \(n \in \mathbf{N}\)、A 中的每个族 \((a_i)_{0 \leqslant i \leqslant n}\) 以及每个复数族 \((t_i)_{0 \leqslant i \leqslant n}\)，有

\[
\sum_ {i = 0} ^ {n} \sum_ {j = 0} ^ {n} \bar {t} _ {i} t _ {j} f (a _ {i}, a _ {j}) = \lambda \left(\left(\sum_ {i = 0} ^ {n} t _ {i} a _ {i}\right) ^ {*} \left(\sum_ {j = 0} ^ {n} t _ {j} a _ {j}\right)\right) \geqslant 0,
\]

从而得到结论（V, p. 433, 定义 1）。

定义 3. — 设 A 是对合巴拿赫代数，且 \(\lambda \in \mathrm{A}_{+}^{\prime}\) 是 A 上的连续正线性泛函。\(\lambda\) 的希尔伯特实现是一个三元组 (E, \(x, \varphi\) )，其中 E 是希尔伯特空间，\(x\) 是 E 的一个元，\(\varphi\) 是从 A 到 \(\mathcal{L}(\mathrm{E})\) 的对合代数态射，并且 \(\lambda(b^{*}a) = \langle \varphi(b)x | \varphi(a)x \rangle\) 对所有 \((a,b) \in \mathrm{A}^{2}\) 都成立。

如果元素 \(\varphi(a)x\)（其中 \(a\in\mathrm{A}\)）在 E 中构成全集，则称 (E, \(x,\varphi\) ) 是 \(\lambda\) 的循环希尔伯特空间实现。

命题 5. --- 设 A 是具有逼近单位元的对合巴拿赫代数。A 上的每个连续正线性泛函都有一个循环希尔伯特实现 (E, x, φ)。若 A 含单位元，则存在一个其中态射 φ 保单位元的此类实现。

由引理 3，b)，不妨设 A 是含单位元的代数。

设 \(\lambda\) 是 A 上的连续正线性泛函。设 (E, \(g\) ) 是 A 上的普遍正核 \(f\) 的循环希尔伯特实现，该核由 \(f(a,b) = \lambda(a^{*}b)\) 定义（命题 4 和定理 1）。映射 \(g\) 连续，其像在 E 中是全的，并且 \(\lambda(a^{*}b) = \langle g(a)|g(b)\rangle\) 对所有 \((a,b) \in \mathrm{A}^2\) 都成立。置 \(x = g(1) \in \mathrm{E}\)。

从 A 到 E 的映射 g 是线性的：事实上，对 \((a,b,c)\in\mathrm{A}^{3}\) 和 \((s,t)\in\mathbf{C}^{2}\)，有

\[
\begin{array}{c} \langle g (c) | g (s a + t b) \rangle = \lambda (c ^ {*} (s a + t b)) \\ = s \lambda (c ^ {*} a) + t \lambda (c ^ {*} b) = \langle g (c) | s g (a) + t g (b) \rangle , \end{array}
\]

由于 \(g\) 的像在 E 中是全的，故断言成立。

特别地，g 的像 F 是 E 的稠密向量子空间。g 的核是 A 的左理想：若 \(g(b) = 0\)，则对 A 中任意 a 和 c，有

\[
\langle g (a b) \mid g (c) \rangle = \lambda ((a b) ^ {*} c) = \lambda (b ^ {*} (a ^ {*} c)) = \langle g (b) \mid g (a ^ {*} c) \rangle = 0,
\]

由于 F 在 E 中稠密，所以 \(g(ab) = 0\)。

设 \(a \in \mathrm{A}\)。由于 \(g\) 的核是 \(\mathrm{A}\) 的左理想，存在一个线性映射 \(\widetilde{\varphi}(a) \colon \mathrm{F} \to \mathrm{F}\)，使得 \(\widetilde{\varphi}(a)(g(b)) = g(ab)\) 对所有 \(b \in \mathrm{A}\) 成立。此外，对每个 \(b \in \mathrm{A}\)，有

\[
\| \widetilde {\varphi} (a) (g (b)) \| ^ {2} = \| g (a b) \| ^ {2} = \lambda \left(b ^ {*} a ^ {*} a b\right) \leqslant \| a ^ {*} a \| \lambda \left(b ^ {*} b\right) \leqslant \| a \| ^ {2} \| g (b) \| ^ {2}
\]

（引理 3，c)），所以 \(\widetilde{\varphi}(a)\) 连续且范数 \(\leqslant \|a\|\)。因此存在唯一的自同态 \(\varphi(a) \in \mathcal{L}(\mathrm{E})\)，它通过对子空间的传递诱导出 \(\widetilde{\varphi}(a)\)。

设 \(a\) 和 \(b\) 属于 A。有

\[
(\varphi (a) \circ \varphi (b)) (g (c)) = g (a b c) = \varphi (a b) (g (c))
\]

对 A 中任意 c 都成立，因此由于 F 在 E 中稠密，\(\varphi(ab) = \varphi(a) \circ \varphi(b)\)。

映射 \(\varphi\) 从 A 到 \(\mathcal{L}(\mathrm{E})\) 是线性的：事实上，对任意 \((a,b)\in\mathrm{A}^{2}\) 和 \((s,t)\in\mathbf{C}^{2}\)，有

\[
\varphi (s a + t b) (g (c)) = g ((s a + t b) c) = (s \varphi (a) + t \varphi (b)) g (c)
\]

对 A 中任意 \(c \in A\) 都成立，从而由于 F 在 E 中稠密，\(\varphi(sa + tb) = s\varphi(a) + t\varphi(b)\)。由于 \(\|\varphi(b)\| \leqslant \|b\|\) 对所有 \(b \in A\) 成立，映射 \(\varphi\) 连续。又有 \(\varphi(1) = 1_{\mathrm{E}}\)，因为 \(\varphi(1)(g(a)) = g(a)\) 对所有 \(a \in A\) 成立，且 F 在 E 中稠密。

最后，设 a、b、c 属于 A。有

\[
\begin{array}{c} \langle g (c) \mid \varphi (a ^ {*}) (g (b)) \rangle = \langle g (c) \mid g (a ^ {*} b) \rangle = \lambda (c ^ {*} a ^ {*} b) \\ = \langle g (a c) \mid g (b) \rangle = \langle \varphi (a) (g (c)) \mid g (b) \rangle \end{array}
\]

由于 F 在 E 中稠密，从而 \(\varphi(a^{*})=\varphi(a)^{*}\)。

总之，映射 \(\varphi\) 是从 A 到 \(\mathcal{L}(\mathrm{E})\) 的连续对合代数态射，且 \(g(a) = \varphi(a)x\) 对所有 \(a \in \mathrm{A}\) 成立；因此三元组 \((\mathrm{E}, x, \varphi)\) 是 \(\lambda\) 的循环希尔伯特实现。

命题 6. --- 设 A 是对合巴拿赫代数，λ 是 A 上的连续正线性泛函，且 (E₁, x₁, φ₁) 和 (E₂, x₂, φ₂) 是 λ 的希尔伯特实现。设 (E₁, x₁, φ₁) 是循环希尔伯特实现。

\begin{enumerate}
\def\labelenumi{\alph{enumi})}
\item
  存在唯一的从 \(E_{1}\) 到 \(E_{2}\) 的连续线性映射 u，它是从 \(\varphi_{1}\) 到 \(\varphi_{2}\) 的表示态射（I, p. 11），并且满足 \(u(\varphi_{1}(a)x_{1}) = \varphi_{2}(a)x_{2}\) 对所有 \(a \in A\) 成立；
\item
  线性映射 u 是等距的；
\item
  若 \(\left(\mathrm{E}_{2},x_{2},\varphi_{2}\right)\) 是循环的，则 u 是同构；
\item
  若 \(\left(\mathrm{E}_2,x_2,\varphi_2\right)\) 是循环的且 A 含单位元，则 \(u(x_{1}) = x_{2}\)。
\end{enumerate}

对 \(j \in \{1,2\}\)，定义 \(\gamma_j \colon A \to E_j\)，使得 \(\gamma_j(a) = \varphi_j(a)x_j\) 对所有 \(a \in A\) 成立。按定义，\((E_j, \gamma_j)\) 是由普遍正核 \(f\) 在 \(A\) 上定义的希尔伯特实现，该核由 \((a,b) \mapsto \lambda(a^*b)\) 定义。希尔伯特实现 \((E_1, \gamma_1)\) 是循环的；由 V, p. 434, 命题 1，存在唯一的连续线性映射 \(u \colon E_1 \to E_2\)，使得 \(\gamma_2 = u \circ \gamma_1\)，且该映射是等距的。此外，若 \((E_2, x_2, \varphi_2)\) 也是循环的，则 \(u\) 是同构。要证明 \(a)\)、\(b)\) 和 \(c)\)，只需证明 \(u\) 是从 \(\varphi_1\) 到 \(\varphi_2\) 的表示态射。

设 \(a \in \mathrm{A}\)。对每个 \(b \in \mathrm{A}\)，有

\[
\begin{array}{c} (u \circ \varphi_ {1} (a)) (\gamma_ {1} (b)) = (u \circ \varphi_ {1} (a b)) x _ {1} = (u \circ \gamma_ {1}) (a b) \\ = \gamma_ {2} (a b) = \varphi_ {2} (a) (\gamma_ {2} (b)) = \varphi_ {2} (a) (u (\gamma_ {1} (b))), \end{array}
\]

因此，连续线性映射 \(u \circ \varphi_1(a)\) 和 \(\varphi_2(a) \circ u\) 在 \(E_1\) 中由 \(\gamma_1\) 的像生成的子空间上一致；由假设，该子空间在 \(E_1\) 中稠密，所以 \(u \circ \varphi_1(a) = \varphi_2(a) \circ u\)。

最后证明 \(d\) )。因此设 A 含单位元，且 \((\mathrm{E}_2, x_2, \varphi_2)\) 是循环的。存在循环希尔伯特空间实现 \((\mathrm{E}_3, x_3, \varphi_3)\) 的 \(\lambda\)，使得 \(\varphi_3\) 是保单位元的态射（命题 5）。设 \(j \in \{1, 2\}\)。将前述结果应用于 \((\mathrm{E}_j, x_j, \varphi_j)\) 和 \((\mathrm{E}_3, x_3, \varphi_3)\)，可知存在等距同构 \(\widetilde{u}_j\)，从 \(\mathrm{E}_j\) 到 \(\mathrm{E}_3\)，使得 \(\widetilde{u}_j \circ \varphi_j(a) = \varphi_3(a) \circ \widetilde{u}_j\) 对所有 \(a \in \mathrm{A}\) 成立。取 \(a = 1_{\mathrm{A}}\)，可知 \(\varphi_j\) 保单位元。于是 \(x_2 = \varphi_2(1_{\mathrm{A}})x_2 = u(\varphi_1(1_{\mathrm{A}})x_1) = u(x_1)\)。

注。--- 保持命题中的记号。可能 \(u(x_{1})\) 不同于 \(x_{2}\)（V, p. 493, 习题 3，b)）。不过，三元组 \((\mathrm{E}_{2}, u(x_{1}), \varphi_{2})\) 也是 \(\lambda\) 的希尔伯特实现。

\subsection{星代数的表示}

设 A 是星代数。以 \(A_{+}\) 表示 A 的正元组成的闭凸锥（I, p. 115, 定义 6）。A 上的线性泛函 \(\lambda\) 是正的，当且仅当 \(\lambda(\mathrm{A}_{+}) \subset \mathbf{R}_{+}\)（I, p. 118, 定理 2）。

命题 7. --- 设 A 是星代数。A 上的每个正线性泛函 \(\lambda\) 都是连续的。

先证明 \(\lambda\) 在 \(\mathrm{A}_{+}\) 与 A 的单位球的交集上有界。反之，存在序列 \((x_{n})_{n\geqslant 1}\) 属于 \(\mathrm{A}_{+}\)，满足 \(\| x_n\| \leqslant 1\) 且 \(\lambda (x_n)\geqslant n\)，对每个整数 \(n\geqslant 1\) 都成立。通项为 \(n^{-2}x_{n}\) 的级数收敛到 A 中的一个元 \(x\)（参见 TG, IV, p. 33, 例 3）。对每个整数 \(\mathrm{N}\geqslant 1\)，由于

\[
\sum_ {n = \mathrm{N} + 1} ^ {+ \infty} \frac {1}{n ^ {2}} x _ {n} \in \mathrm{A} _ {+},
\]

由（I, p. 116, 命题 14）可得

\[
\sum_ {n = 1} ^ {\mathrm{N}} \frac {1}{n} \leqslant \sum_ {n = 1} ^ {\mathrm{N}} \frac {1}{n ^ {2}} \lambda (x _ {n}) = \lambda (x) - \lambda \left(\sum_ {n = \mathrm{N} + 1} ^ {+ \infty} \frac {1}{n ^ {2}} x _ {n}\right) \leqslant \lambda (x),
\]

这与事实矛盾（TG, IV, p. 33, 例 4）。

由于 A 的单位球中的每个元都是 A 中至多四个范数 \(\leqslant 1\) 的正元的系数绝对值不超过 1 的线性组合（I, p. 96, 引理 2 及 I, p. 117, 公式 (4)），所以 \(\lambda\) 连续。

存在一些对合巴拿赫代数，它们具有不连续的正线性泛函（V, p. 493, 习题 3，a)）。

命题 8. --- 设 A 是星代数，a 是 A 的非零元。存在正线性泛函 \(\lambda \in \mathrm{A}_{+}^{\prime}\)，使得 \(\lambda(a^{*}a) > 0\)。

考虑 A 的埃尔米特元组成的实巴拿赫空间 \(A_{h}\)。集合 \(A_{+}\) 是 \(A_{h}\) 中的闭凸尖锥（I, p. 116, 命题 14）。元 \(a^{*}a\) 是正的（I, p. 118, 定理 2）且非零，故埃尔米特元 \(-a^{*}a\) 不是正的。由 EVT, II, p. 42, 推论 5，存在连续实线性泛函 \(\lambda\in A_{h}^{\prime}\)，使得 \(\lambda(-a^{*}a)<0\) 且 \(\lambda(\mathrm{A}_{+})\subset\mathbf{R}_{+}\)。线性泛函 \(\lambda\) 延拓为 A 上的埃尔米特 C-线性泛函（参见 I, p. 98），该泛函是正的，并具有所要求的性质。

定理 2（盖尔范德--奈马克）。--- 设 A 是星代数。存在希尔伯特空间 E 以及 \(\varphi\) 从 A 到 \(\mathcal{L}(\mathrm{E})\) 的等距对合代数态射。若 A 含单位元，则存在这样的保单位元态射。

对每个 \(b \in \mathrm{A} - \{0\}\)，取 A 上的连续正线性泛函 \(\lambda_b\)，使得 \(\lambda_b(b^*b) > 0\)（命题 8）。由于 A 具有逼近单位元（I, p. 121, 命题 18），存在希尔伯特实现 \((\mathrm{E}_b, x_b, \varphi_b)\) 的 \(\lambda_b\)（命题 5）。若 A 含单位元，可设 \(\varphi_b\) 保单位元（同上）。有 \(\| \varphi_b(b) \|^2 = \lambda_b(b^*b) \neq 0\)。

令 E 为各空间 \(E_{b}\)（其中 \(b\) 属于 \(A - \{0\}\)）的外希尔伯特直和。对每个 \(a \in A\)，以 \(\varphi(a) \in \mathcal{L}(E)\) 表示唯一的连续线性映射，其在每个 \(E_{b}\) 上的限制都与 \(\varphi_{b}(a)\) 一致，其中 \(b \in A - \{0\}\)。映射 \(a \mapsto \varphi(a)\) 是对合代数态射；它是单射，因而是等距的（I, p. 112, 命题 9）；若 A 含单位元，还满足 \(\varphi(1) = 1_{\mathrm{E}}\)。证明完毕。

*因此，以星代数为对象、以对合代数态射为态射的范畴，等价于以希尔伯特空间的自同态代数的闭对合子代数为对象、以对合代数态射为态射的范畴。*

\subsection{拓扑群上的正定函数}

在本节中，G 是拓扑群，其单位元记为 e。所考虑的希尔伯特空间都是复的。

定义 4. --- 连续函数 \(\varphi\in\mathcal{C}(\mathrm{G})\) 称为 G 上的正定函数，如果函数 \(f\) 由 \(f(g,h)=\varphi(g^{-1}h)\) 定义在 \(\mathrm{G}\times\mathrm{G}\) 上，且是 \(\mathrm{G}\) 上的普遍正核。

以 \(\operatorname{Pos}(\mathbf{G})\) 表示 G 上正定函数的集合，并以 \(\operatorname{Pos}_{1}(\mathbf{G})\)（分别以 \(\operatorname{Pos}_{\leqslant1}(\mathbf{G})\)）表示 \(\varphi\in\operatorname{Pos}(\mathbf{G})\) 中满足 \(\varphi(e)=1\)（分别满足 \(\varphi(e)\leqslant1\)）的函数所成的子集。

例。--- 设 \(\varrho\) 是 G 在希尔伯特空间 E 上的酉表示，且 \(x \in \mathrm{E}\)。设 \(\varphi\) 是由 \(\varphi(g) = \langle x \mid \varrho(g)x \rangle\) 对所有 \(g \in \mathrm{G}\) 定义的对角矩阵系数；函数 \(\varphi\) 连续。由此可得

\[
\varphi (g ^ {- 1} h) = \langle x | \varrho (g) ^ {*} \varrho (h) x \rangle = \langle \varrho (g) x | \varrho (h) x \rangle
\]

对所有 \((g,h)\in\mathrm{G}\times\mathrm{G}\) 都成立。因此 \(\varphi\in\mathrm{Pos}(\mathrm{G})\)（V, p. 432, 定理 1, (iv)）。\(\varphi\in\mathrm{Pos}_{1}(\mathrm{G})\) 当且仅当 \(\|x\|=1\)。

定义 5. --- 设 \(\varphi\) 是 G 上的正定函数。\(\varphi\) 的希尔伯特实现是一个对 \((\varrho, x)\)，其中 \(\varrho\) 是 G 在希尔伯特空间 E 上的酉表示，\(x \in \mathrm{E}\)，且 \(\varphi(g) = \langle x \mid \varrho(g)x \rangle\) 对所有 \(g \in \mathrm{G}\) 成立。

若 x 是 \(\varrho\) 的循环向量，则称其为 \(\varphi\) 的循环希尔伯特实现。

命题 9. --- 设 \(\varphi\in\mathcal{C}(\mathrm{G})\)。以下条件等价：

\begin{enumerate}
\def\labelenumi{(\roman{enumi})}
\item
  函数 \(\varphi\) 是 G 上的正定函数；
\item
  存在 \(\varphi\) 的循环希尔伯特实现；
\item
  存在 \(\varphi\) 的希尔伯特实现。
\end{enumerate}

条件 (ii) 蕴含条件 (iii)，而由上述例子，条件 (iii) 蕴含条件 (i)。

最后证明 (i) 蕴含 (ii)。设 (E, \(\gamma\) ) 是由 \(f(g,h)=\varphi(g^{-1}h)\) 定义的普遍正核（其中 \((g,h)\in\mathrm{G}\times\mathrm{G}\)）的循环希尔伯特实现。设 \(k\in\mathrm{G}\)。G 上的连续函数 \(\gamma_{k}:g\mapsto\gamma(kg)\) 满足

\[
\langle \gamma_ {k} (g) | \gamma_ {k} (h) \rangle = \langle \gamma (k g) | \gamma (k h) \rangle = f (k g, k h) = \varphi ((k h) ^ {- 1} k g) = f (g, h)
\]

对所有 \((g,h)\in\mathrm{G}\times\mathrm{G}\) 都成立，因此 \((\mathrm{E},\gamma_{k})\) 是 f 的希尔伯特实现。由 V, p. 434, 命题 1，存在唯一的酉算子 \(\varrho(k)\) 属于 \(\mathcal{L}(\mathrm{E})\)，使得 \(\gamma_{k}=\varrho(k)\circ\gamma\)。对每个 \(g\in G\) 以及每个 \((k_{1},k_{2})\in\mathrm{G}\times\mathrm{G}\)，可得

\[
\varrho (k _ {1} k _ {2}) (\gamma (g)) = \gamma (k _ {1} k _ {2} g) = \varrho (k _ {1}) (\varrho (k _ {2}) (\gamma (g))),
\]

从而 \(\varrho(k_1k_2) = \varrho(k_1)\varrho(k_2)\)，因为 \(\gamma\) 的像是 E 的全集。

设 \(k \in G\)。对每个 \(g \in G\)，有 \(\varrho(g)(\gamma(k)) = \gamma(gk)\)，因此由 \(g \mapsto \varrho(g)(\gamma(k))\) 定义的从 G 到 E 的映射连续。由于自同态 \(\varrho(g)\) 是酉的，故可推出 \(\varrho\) 是 G 在 E 上的酉表示（V, p. 380, 引理 4）。再置 \(x = \gamma(e)\)。向量集 \(\varrho(g)x = \varrho(g)(\gamma(e)) = \gamma(g)\) 在 E 中是全的，其中 \(g\) 遍历 G，故 x 是 \(\varrho\) 的循环向量。由于

\[
\langle x \mid \varrho (g) x \rangle = f (e, g) = \varphi (g)
\]

对所有 \(g \in G\) 都成立，所以对 \((\varrho, x)\) 是 \(\varphi\) 的循环希尔伯特实现。

命题 10. — 设 \(\varphi\) 是 G 上的正定函数，\((\varrho_{1}, x_{1})\) 和 \((\varrho_{2}, x_{2})\) 是 \(\varphi\) 的希尔伯特实现，且希尔伯特表示 \((\varrho_{1}, x_{1})\) 是循环的。存在唯一的等距 G-态射 \(u\)，从 \(\varrho_{1}\) 到 \(\varrho_{2}\)，使得 \(u(x_{1}) = x_{2}\)。若 \((\varrho_{2}, x_{2})\) 也是循环的，则 \(u\) 是同构。

对 \(1 \leqslant j \leqslant 2\)，令 \(\mathrm{E}_j\) 为 \(\varrho_j\) 的空间，并令 \(\gamma_j\) 为 \(\mathbf{G}\) 上由 \(\gamma_j(g) = \varrho_j(g)x_j\) 对所有 \(g \in \mathbf{G}\) 定义的函数。对 \((\mathrm{E}_1, \gamma_1)\) 和 \((\mathrm{E}_2, \gamma_2)\) 是正定函数 \((g, h) \mapsto \varphi(g^{-1}h)\) 的希尔伯特实现，且 \((\mathrm{E}_1, \gamma_1)\) 是循环希尔伯特实现。由 V, p. 434, 命题 1，存在唯一的等距线性映射 \(u: \mathrm{E}_1 \to \mathrm{E}_2\)，使得 \(\gamma_2 = u \circ \gamma_1\)。特别地，有 \(x_2 = \gamma_2(e) = u(\gamma_1(e)) = u(x_1)\)。此外，对任意 \(g\) 和 \(h\) 属于 \(\mathbf{G}\)，有

\[
\varrho_ {2} (g) u \left(\gamma_ {1} (h)\right) = \varrho_ {2} (g) \gamma_ {2} (h) = \gamma_ {2} (g h) = u \left(\gamma_ {1} (g h)\right) = u \left(\varrho_ {1} (g) \gamma_ {1} (h)\right),
\]

由于元素 \(\gamma_{1}(h)\)（其中 \(h \in G\)）所成的集合在 \(E_{1}\) 中是全的，这意味着 u 是 G-态射。根据同上，若 \((\varrho_{2}, x_{2})\) 也是循环的，则 u 是等距同构。

命题 11. — 设 \(\varphi \in \mathrm{Pos}_{1}(\mathrm{G})\)，且 \((\varrho, x)\) 是 \(\varphi\) 的循环希尔伯特实现。则 \(\varrho\) 是 G 的不可约表示，当且仅当 \(\varphi\) 是 \(\mathrm{Pos}_{1}(\mathrm{G})\) 的极点（EVT, II, p. 57, 定义 1）。

令 E 为 \(\varrho\) 的空间。先设 \(\varrho\) 不是不可约的。令 F 是 E 的一个在 \(\varrho\) 下稳定的、非零且不同于 E 的闭子空间。写作 \(x = x_1 + x_2\)，其中 \(x_1 \in \mathrm{F}\)，\(x_2 \in \mathrm{F}^\circ\)。于是 \(1 = \|x\|^2 = \|x_1\|^2 + \|x_2\|^2\)。由 \(x_1\) 生成的 E 的子表示包含于 F 中，所以 \(x_1 \neq x\)，因为 \(x\) 是 \(\varrho\) 的循环向量，从而 \(x_2 \neq 0\)。同理可验证 \(x_1 \neq 0\)。

对 \(j = 1\) 和 \(j = 2\)，令 \(\varphi_j\) 为 G 上满足

\[
\varphi_ {j} (g) = \frac {1}{\| x _ {j} \| ^ {2}} \langle x _ {j} | \varrho (g) x _ {j} \rangle
\]

对所有 \(g \in G\) 均成立。我们有 \(\varphi_{j} \in \mathrm{Pos}_{1}(G)\)（V, p. 443, 例）。由于 \(\varphi = \|x_{1}\|^{2}\varphi_{1} + \|x_{2}\|^{2}\varphi_{2}\)，只需验证 \(\varphi_{1} \neq \varphi_{2}\) 即可证明 \(\varphi\) 不是 \(\mathrm{Pos}_{1}(G)\) 的极点；为此只需证明 \(\varphi \neq \varphi_{1}\)。

反证。假设 \(\varphi = \varphi_1\)。由于 \(\langle x_1 | \varrho(g)x_2 \rangle = 0\) 对所有 \(g \in G\) 成立，应有

\[
\frac {1}{\| x _ {1} \| ^ {2}} \langle x _ {1} | \varrho (g) x \rangle = \frac {1}{\| x _ {1} \| ^ {2}} \langle x _ {1} | \varrho (g) x _ {1} \rangle = \varphi_ {1} (g) = \varphi (g) = \langle x | \varrho (g) x \rangle
\]

对所有 \(g \in G\) 都成立，从而 E 的向量子空间中的每个元素 y 都满足 \(\langle x_{1} | y \rangle = \langle \|x_{1}\|^{2}x | y \rangle\)，其中该子空间由元素 \(\varrho(g)x\)（其中 \(g \in G\)）生成；因此对于所有 \(y \in E\)，由于 x 是 \(\varrho\) 的循环向量，这都成立。这将意味着 \(x_{1} = \|x_{1}\|^{2}x\) 也是 \(\varrho\) 的循环向量，矛盾，故断言成立。

现在设 \(\varrho\) 不可约，并证明 \(\varphi\) 是 \(\mathrm{Pos}_1(\mathrm{G})\) 的极点。设 \(\varphi_1 \neq \varphi_2\) 是 \(\mathrm{Pos}_1(\mathrm{G})\) 中的元，且 \(t_1, t_2 \in [0,1]\) 满足 \(t_1 + t_2 = 1\) 和 \(\varphi = t_1\varphi_1 + t_2\varphi_2\)。对 \(j \in \{1,2\}\)，设 \((\varrho_j, x_j)\) 是 \(\varphi_j\) 的循环希尔伯特实现，并令 \(\mathrm{E}_j\) 为 \(\varrho_j\) 的空间。令 \(x_3 = t_1^{1/2}x_1 + t_2^{1/2}x_2\)。则 \((\varrho_1 \oplus \varrho_2, x_3)\) 是 \(\varphi\) 的希尔伯特实现。由于 \((\varrho, x)\) 是循环的，存在等距 G-态射 \(u: \mathrm{E} \to \mathrm{E}_1 \oplus \mathrm{E}_2\)，使得 \(u(x) = x_3\)（命题 10）。

令 \(j = 1\) 或 \(j = 2\)。由于 \(\varrho\) 不可约，存在 \(\lambda_j \geqslant 0\)，使得 G-态射 \(u_j = \mathrm{pr}_j \circ u\) 从 E 到 \(E_j\) 满足

\[
\langle u _ {j} (y) \mid u _ {j} (y ^ {\prime}) \rangle = \lambda_ {j} \langle y \mid y ^ {\prime} \rangle
\]

对 E 中任意 \(y\) 和 \(y'\) 都成立（若 \(u_j \neq 0\)，这是 V, p. 388, 推论 5；否则可取 \(\lambda_j = 0\)）。对每个 \(g \in G\)，可得

\[
\begin{array}{r} t _ {j} \varphi_ {j} (g) = \langle t _ {j} ^ {1 / 2} x _ {j} | \varrho_ {j} (g) (t _ {j} ^ {1 / 2} x _ {j}) \rangle = \langle u _ {j} (x) | \varrho_ {j} (g) (u _ {j} (x)) \rangle \\ = \langle u _ {j} (x) | u _ {j} (\varrho (g) x) \rangle = \lambda_ {j} \varphi (g). \end{array}
\]

由于 \(\varphi_{j}(e) = \varphi(e) = 1\)，可推出 \(\varphi_{j} = \varphi\) 当 \(t_{j} \neq 0\) 时成立。因此假设 \(\varphi_{1} \neq \varphi_{2}\) 蕴含 \(t_{1}\) 或 \(t_{2}\) 必为零，这就证明了 \(\varphi\) 是 \(\mathrm{Pos}_{1}(\mathrm{G})\) 的极点。

引理 4. --- 设 \(\varphi\in\mathrm{Pos}(\mathbf{G})\)。函数 \(\varphi\) 在 \(\mathbf{G}\) 上有界，且 \(\|\varphi\|_{\infty}=\varphi(e)\)。此外，有

\[
| \varphi (g ^ {- 1} h) - \varphi (h) | \leqslant \sqrt {2 \varphi (e) (\varphi (e) - \mathcal {R} (\varphi (g)))}\tag{1}
\]

对所有 \((g,h)\in \mathbf{G}\times \mathbf{G}\) 都成立。

设 \((\varrho, x)\) 是 \(\varphi\) 的希尔伯特实现。有 \(\varphi(e) = \|x\|^2\)。因此对每个 \(g \in G\)，由柯西–施瓦茨不等式可得 \(|\varphi(g)| = |\langle x | \varrho(g)x\rangle| \leqslant \varphi(e)\)。这就证明了第一个断言。

对每个 \((g,h)\in \mathrm{G}\times \mathrm{G}\)，有

\[
\varphi (g ^ {- 1} h) - \varphi (h) = \langle x | \varrho (g ^ {- 1} h) x \rangle - \langle x | \varrho (h) x \rangle = \langle \varrho (g) x - x | \varrho (h) x \rangle
\]

由于 \(\varrho\) 是酉的，从而

\[
| \varphi (g ^ {- 1} h) - \varphi (h) | \leqslant \| \varrho (g) x - x \| \| \varrho (h) x \| \leqslant \varphi (e) ^ {1 / 2} \| \varrho (g) x - x \|.
\]

注意到

\[
\| \varrho (g) x - x \| ^ {2} = 2 \| x \| ^ {2} - 2 \mathcal {R} (\langle x | \varrho (g) x \rangle) = 2 (\varphi (e) - \mathcal {R} (\varphi (g))).
\]

注。--- 集合 \(\operatorname{Pos}(G)\)（分别为 \(\operatorname{Pos}_{1}(G)\) 和 \(\operatorname{Pos}_{\leqslant1}(G)\)）是星代数 \(\mathcal{C}_{b}(G)\) 的凸自伴子集。它们在赋予逐点收敛拓扑的空间 \(\mathcal{C}_{b}(G)\) 中是闭的。集合 \(\operatorname{Pos}(G)\) 是实巴拿赫空间 \(\mathcal{C}_{b}(G)\) 中顶点为 0 的凸锥。

\subsection{局部紧群的酉对偶}

设 G 是局部紧拓扑群。给 G 赋予左哈尔测度 \(\mu\)。对 \(p \in [1, +\infty]\)，以 \(\mathcal{L}^{p}(\mathrm{G})\)（分别以 \(\mathrm{L}^{p}(\mathrm{G})\)）表示空间 \(\mathcal{L}_{\mathbf{C}}^{p}(\mathrm{G}, \mu)\)（分别表示空间 \(\mathrm{L}_{\mathbf{C}}^{p}(\mathrm{G}, \mu)\)）。将空间 \(\mathcal{C}_{b}(\mathrm{G})\) 与其在 \(\mathrm{L}^{\infty}(\mathrm{G})\) 中的像等同起来。

以 \(\Delta\) 表示 G 的模函数。回忆 \(L^1(G)\) 是对合巴拿赫代数，其对合由映射 \(f \mapsto f^*\) 通过取商诱导，其中 \(f^*(y) = \Delta^{-1}(y)\overline{f(y^{-1})}\) 对所有 \(f \in \mathcal{L}^1(G)\) 和 \(y \in G\) 成立（参见 I, p. 99, 例 4）。根据 INT, VIII, p. 172, §4, no. 7, 命题 20，巴拿赫代数 \(L^1(G)\) 具有逼近单位元。

设 \(\varrho\) 是 G 在复希尔伯特空间 E 上的酉表示。映射 \(f\mapsto\varrho(f)\) 是从 L \(^{1}\) (G) 到 \(\mathcal{L}\) (E) 的连续对合代数态射（V, p. 401, 引理 1）；它是 L \(^{1}\) (G) 在 E 上的非退化表示（INT, VIII, p. 139, § 2, no. 7, 命题 10, (i)）。

命题 12. --- 设 E 是复希尔伯特空间，\(\widetilde{\pi}\) 是从 L \(^{1}\) (G) 到 \(\mathcal{L}\) (E) 的对合代数态射。若表示 \(\widetilde{\pi}\) 非退化，则存在 G 在 E 上的唯一酉表示 \(\pi\)，使得 \(\widetilde{\pi}(f)=\pi(f)\) 对所有 \(f\in\mathrm{L}^{1}(\mathrm{G})\) 成立。

唯一性由 INT, VIII, p. 139, § 2, n° 7, 引理 4 的推论 3 得出。

证明 \(\pi\) 的存在性。由 I, p. 104, 命题 2，有 \(\| \widetilde{\pi}\| \leqslant 1\)。令 F 为 E 中由形如 \(\widetilde{\pi}(f)x\) 的向量生成的子空间，其中 \(f\) 属于 \(\mathrm{L}^1 (\mathrm{G})\)，\(x\) 属于 E；由于表示 \(\widetilde{\pi}\) 非退化，F 在 E 中稠密。

对 e 的每个紧邻域 V，令 \(\varphi_{V}\) 为支集包含于 V 且满足 \(\int\varphi_{V}\mu=1\) 的连续正函数。令 B 为 e 的紧邻域滤子的一个基。

设 \(g\) 属于 \(G\)，设 \(f\) 属于 \(L^1(G)\)。有 \((\varphi_V * \varepsilon_g) * f \to \varepsilon_g * f\) 在 \(L^1(G)\) 中沿 \(\mathfrak{B}\) 的截面滤子收敛（INT, VIII, p. 172, § 4, no. 7, 命题 20），所以 \(\widetilde{\pi}(\varphi_V * \varepsilon_g)\widetilde{\pi}(f)\) 收敛到 \(\widetilde{\pi}(\varepsilon_g * f)\)（在 \(\mathcal{L}(E)\) 中）。这意味着 \(\widetilde{\pi}(\varphi_V * \varepsilon_g)\) 沿 \(\mathfrak{B}\) 的截面滤子逐点收敛于映射 \(\pi(g)\)，从 \(F\) 到 \(F\)；由于 \(\| \widetilde{\pi}(\varphi_V * \varepsilon_g) \| \leqslant 1\) 对每个 \(V\)（它是 \(e\) 的紧邻域）成立，该映射由 EVT, III, p. 26, 推论 3 可知是连续的。因此存在唯一的 \(E\) 的自同态，它通过对子空间的传递诱导出 \(\pi(g)\)。仍以 \(\pi(g)\) 表示该自同态；有 \(\| \pi(g) \| \leqslant 1\)。

设 \(f \in \mathrm{L}^1(\mathrm{G})\)。由于 \(\widetilde{\pi}(\varphi_{\mathrm{V}} * \varepsilon_g) \widetilde{\pi}(f)\) 收敛到 \(\widetilde{\pi}(\varepsilon_g * f)\)，有 \(\pi(g) \widetilde{\pi}(f) = \widetilde{\pi}(\varepsilon_g * f)\)。

当 \(g = e\) 时，该关系表明 \(\pi(e)\) 在 \(F\) 上是恒等映射，因而 \(\pi(e) = 1_{E}\)。设 \(g_1\) 和 \(g_2\) 属于 \(G\)。有

\[
\pi (g _ {1}) \pi (g _ {2}) \widetilde {\pi} (f) = \pi (g _ {1}) \widetilde {\pi} (\varepsilon_ {g _ {2}} * f) = \widetilde {\pi} (\varepsilon_ {g _ {1} g _ {2}} * f) = \pi (g _ {1} g _ {2}) \widetilde {\pi} (f),
\]

所以 \(\pi(g_1)\pi(g_2) = \pi(g_1g_2)\) 在 F 上成立，因而在 E 上也成立。这证明了映射 \(g \mapsto \pi(g)\) 是 G 在 E 上的表示。由于 \(\|\pi(g)\| \leqslant 1\) 且 \(\|\pi(g)^{-1}\| = \|\pi(g^{-1})\| \leqslant 1\)，E 的自同态 \(\pi(g)\) 是等距的，因而是酉的（EVT, V, p. 40, 命题 3）。

设 \(f \in \mathrm{L}^{1}(\mathrm{G})\) 且 \(x \in E\)。映射 \(g \mapsto \varepsilon_{g} * f\) 从 G 到 \(\mathrm{L}^{1}(\mathrm{G})\) 连续（INT, VIII, p. 136, § 2, no. 5 及 p. 144, § 3, no. 2, 公式 (5)），且 \(\widetilde{\pi}\) 连续，因此映射 \(g \mapsto \widetilde{\pi}(\varepsilon_{g} * f)x = \pi(g)\widetilde{\pi}(f)x\) 在 \(g = e\) 处连续。由于 F 在 E 中稠密，表示 \(\pi\) 是酉表示（V, p. 380, 引理 4）。

设 \(f_1\) 和 \(f_2\) 属于 \(L^1(G)\)。根据 INT, VIII, p. 127, §1, no. 5, 命题 7，有关系

\[
f _ {1} * f _ {2} = \int_ {\mathrm{G}} f _ {1} (g) (\varepsilon_ {g} * f _ {2}) d \mu (g)
\]

在 L \(^{1}\) (G) 中，从而

\[
\begin{array}{l} \widetilde {\pi} (f _ {1}) \widetilde {\pi} (f _ {2}) = \widetilde {\pi} (f _ {1} * f _ {2}) = \int_ {\mathrm{G}} f _ {1} (g) \widetilde {\pi} (\varepsilon_ {g} * f _ {2}) d \mu (g) \\ \qquad = \int_ {\mathrm{G}} f _ {1} (g) \pi (g) \widetilde {\pi} (f _ {2}) d \mu (g) = \Big (\int_ {\mathrm{G}} f _ {1} (g) \pi (g) d \mu (g) \Big) \widetilde {\pi} (f _ {2}), \end{array}
\]

利用 INT, VI, p. 9, § 1, no. 1, 命题 1，可得

\[
\widetilde {\pi} (f) = \int_ {\mathrm{G}} f (g) \pi (g) d \mu (g) = \pi (f)
\]

对所有 \(f \in \mathrm{L}^{1}(\mathrm{G})\) 都成立。命题得证。

命题 13. --- 设 \(\varphi\in\mathrm{L}^{\infty}(\mathrm{G})\)。以 \(\lambda_{\varphi}\) 表示由 \(f\mapsto\langle f,\varphi\rangle\) 定义在 \(\mathrm{L}^{1}(\mathrm{G})\) 上的连续线性泛函。则 \(\varphi\) 是 G 上某个正定函数的等价类，当且仅当 \(\lambda_{\varphi}\) 是 \(\mathrm{L}^{1}(\mathrm{G})\) 上的连续正线性泛函。

设 \(\lambda_{\varphi}\) 是 \(L^{1}(G)\) 上的连续正线性泛函。令 \((E, x, \widetilde{\pi})\) 为 \(\lambda\) 的循环希尔伯特实现（V, p. 438, 命题 5）。

令 \(\pi\) 为 G 在 E 上的酉表示，使得 \(\pi(f) = \widetilde{\pi}(f)\) 对所有 \(f \in \mathrm{L}^1(\mathrm{G})\) 成立（命题 12）。于是 \(\lambda_{\varphi}(f) = \langle x | \pi(f)x \rangle\) 对所有 \(f \in \mathrm{L}^1(\mathrm{G})\) 成立。

由于

\[
\pi (f) = \int_ {\mathrm{G}} f (g) \pi (g) d \mu (g)
\]

对所有 \(f\in \mathrm{L}^1 (\mathrm{G})\) 都成立，从而

\[
\int_ {\mathrm{G}} f (g) \varphi (g) d \mu (g) = \lambda_ {\varphi} (f) = \langle x | \pi (f) x \rangle = \int_ {\mathrm{G}} f (g) \langle x | \pi (g) x \rangle d \mu (g)
\]

对所有 \(f \in \mathrm{L}^{1}(\mathrm{G})\) 都成立。因此，\(\varphi\) 是 \(\mathrm{L}^{\infty}(\mathrm{G})\) 中由 G 上的函数 \(g \mapsto \langle x | \pi(g)x \rangle\) 所定义的等价类，而该函数是 G 上的正定函数（V, p. 443, 命题 9）。

反之，设 \(\varphi\in\mathrm{Pos}(\mathrm{G})\)。令 \(f\in\mathcal{K}(\mathrm{G})\)。于是有

\[
\begin{array}{c} \lambda_ {\varphi} (f ^ {*} * f) = \int_ {G} \varphi (y) \Big (\int_ {G} \Delta (z) ^ {- 1} \overline {{f (z ^ {- 1})}} f (z ^ {- 1} y) d \mu (z) \Big) d \mu (y) \\ = \int_ {G} \varphi (y) \Big (\int_ {G} \overline {{f (z)}} f (z y) d \mu (z) \Big) d \mu (y) \end{array}
\]

G × G 上由 \((z, y) \mapsto \varphi(y)\overline{f(z)}f(zy)\) 定义的连续函数是有界的，并且

\[
\begin{array}{c} \int_ {\mathrm{G}} ^ {*} | \varphi (y) | \Big (\int_ {\mathrm{G}} ^ {*} | \overline {{f (z)}} f (z y) | d \mu (z) \Big) d \mu (y) \\ \leqslant \varphi (e) \int_ {\mathrm{G}} ^ {*} \int_ {\mathrm{G}} ^ {*} | f (z) | | f (z y) | d \mu (z) d \mu (y) \\ = \varphi (e) \Big (\int_ {\mathrm{G}} ^ {*} | f (z) | d \mu (z) \Big) ^ {2} \end{array}
\]

由 INT, V, p. 94, § 8, no. 3, 命题 8，利用勒贝格–富比尼定理（INT, V, p. 96, § 8, no. 4, 定理 1），可推出

\[
\begin{array}{c} \lambda_ {\varphi} (f ^ {*} * f) = \int_ {\mathrm{G}} \overline {{f (z)}} \Bigl (\int_ {\mathrm{G}} \varphi (y) f (z y) d \mu (y) \Bigr) d \mu (z) \\ = \int_ {\mathrm{G} \times \mathrm{G}} \overline {{f (z)}} f (y) \varphi (z ^ {- 1} y) d (\mu \otimes \mu) (z, y). \end{array}
\]

由 V, p. 432, 定理 1, (i)，可知 \(\lambda_{\varphi}(f^{*} * f) \geqslant 0\)，这是将其应用于 \(\mu\) 在 \(f\) 的支集上诱导的测度的结果（INT, IV, p. 186, § 5, n° 7, 定义 4）。由连续性可推出，\(\lambda_{\varphi}(f^{*} * f) \geqslant 0\) 对所有 \(f \in \mathrm{L}^1(\mathrm{G})\) 成立。

空间 \(\mathrm{L}^{\infty}(\mathrm{G})\) 和 \(\mathcal{C}_{b}(\mathrm{G})\) 赋予其巴拿赫空间拓扑，并以 \(f \mapsto \|f\|_{\infty}\) 表示范数。这里将 \(\mathrm{L}^{\infty}(\mathrm{G})\)、\(\mathcal{C}_{b}(\mathrm{G})\) 或 \(\operatorname{Pos}(\mathrm{G})\) 上由弱拓扑 \(\sigma(\mathrm{L}^{\infty}(\mathrm{G}), \mathrm{L}^{1}(\mathrm{G}))\) 诱导的拓扑称为弱拓扑。

由于 \(L^{\infty}(G)\) 与 \(L^{1}(G)\) 的对偶空间相等同（INT, V, p. 61, §5, no. 8, 定理 4），\(L^{\infty}(G)\) 的每个闭球对于弱拓扑都是紧的（EVT, III, p. 17, 推论 3）。

推论。--- 在赋予弱拓扑的 \(\mathrm{L}^{\infty}(\mathrm{G})\) 中，集合 \(\operatorname{Pos}(\mathrm{G})\) 是闭的，集合 \(\operatorname{Pos}_{\leqslant 1}(\mathrm{G})\) 是紧的。

第一个断言由命题得出，第二个断言则由 EVT, III, 同上得出。

一般而言，\(\mathrm{Pos}_{1}(\mathrm{G})\) 对弱拓扑并不紧，群 R 的例子说明了这一点（V, p. 497, 习题 10）。

命题 14. --- \(\mathrm{Pos}_{\leqslant 1}(\mathrm{G})\) 的极点集合等于 \(\mathrm{Pos}_1(\mathrm{G})\) 的极点集合与零函数的并。此外，\(\mathrm{Pos}_1(\mathrm{G})\) 的极点集合的闭凸包包含 \(\mathrm{Pos}_1(\mathrm{G})\)。

由命题 13 的推论，集合 \(\mathrm{Pos}_{\leqslant 1}(\mathrm{G})\) 是尖凸锥 \(\mathrm{Pos}(\mathrm{G})\) 的一个截面，该截面由规范 \(\varphi \mapsto \varphi(e)\) 定义（EVT, II, p. 61, 定义 3 及命题 4）。因此，它的极点是零函数以及元素 \(\varphi\)，它属于 \(\mathrm{Pos}_1(\mathrm{G})\) 且属于 \(\mathrm{Pos}(\mathrm{G})\) 的极端生成元（EVT, II, p. 62, 推论 1）。这些正是 \(\mathrm{Pos}_1(\mathrm{G})\) 的极点（EVT, II, p. 61, 命题 3）。第一个断言由此得出。

证明第二个断言。设 \(\varphi\in\mathrm{Pos}_{1}(\mathrm{G})\)。由于集合 \(\mathrm{Pos}_{\leqslant 1}(\mathrm{G})\) 对弱拓扑是紧的（命题 13 的推论），存在滤子 \(\mathfrak{F}\)，在 \(\mathrm{Pos}_{\leqslant 1}(\mathrm{G})\) 的极点集合的凸包上收敛于 \(\varphi\)（EVT, II, p. 59, 定理 1）。由于巴拿赫空间 \(\mathrm{L}^{\infty}(\mathrm{G})\) 的闭球对于弱拓扑是闭的，且 \(\|\varphi\|_{\infty}=1\)，有 \(\lim_{\psi,\mathfrak{F}}\|\psi\|_{\infty}=1\)（事实上，否则将存在实数 \(c<1\) 以及滤子 \(\mathfrak{G}\)，它在半径为 \(c\) 的 \(\mathrm{L}^{\infty}(\mathrm{G})\) 闭球上，且比 \(\mathfrak{F}\) 更细，这将意味着 \(\varphi\) 属于该闭球）。

根据对 \(\mathrm{Pos}_{\leqslant1}(\mathrm{G})\) 极点的描述，极点的凸包中的每个元 \(\psi\) 都是 G 上具有如下形式的 \(\mathrm{Pos}_{\leqslant1}(\mathrm{G})\) 上的正定函数

\[
\psi = \sum_ {i \in \mathrm{I}} t _ {i} \psi_ {i},
\]

其中 I 是有限集，\(\psi_{i}\) 是 \(\mathrm{Pos}_{1}(\mathrm{G})\) 的极点，对所有 \(i \in I\)，且 \(t_{i} \in [0,1]\)。有 \(\sum_{i} t_{i} = \psi(e) = \|\psi\|_{\infty} \leqslant 1\)（V, p. 446, 引理 4）。若 \(\psi \neq 0\)，函数

\[
\frac {\psi}{\| \psi \| _ {\infty}} = \sum_ {i \in \mathrm{I}} \frac {t _ {i}}{\psi (e)} \psi_ {i}
\]

因此属于 \(\mathrm{Pos}_{1}(\mathrm{G})\) 极点的凸包。由于 \(\lim_{\psi,\mathfrak{F}}\|\psi\|_{\infty}^{-1}\psi=\varphi\)，可知 \(\varphi\) 属于 \(\mathrm{Pos}_{1}(\mathrm{G})\) 极点的闭凸包。断言得证。

\subsection{不可约表示的存在性}

沿用上一节的记号。

定理 3（赖科夫）。--- Pos1(G) 上的弱拓扑与紧收敛拓扑重合。

先证明几个引理。

引理 5. --- 设 X 是局部紧拓扑空间，\(\nu\) 是 X 上的正测度。在 \(\mathcal{C}_{b}(X)\) 的每个有界子集上，由弱拓扑 \(\sigma(L^{\infty}(X), L^{1}(X))\) 诱导的拓扑比紧收敛拓扑粗。

设 B 是 \(\mathcal{C}_{b}(\mathrm{X})\) 的有界子集，且 \(M \in R_{+}\) 满足 \(\|\varphi\|_{\infty} \leqslant M\) 对所有 \(\varphi \in B\)。设 \(\psi \in L^{1}(\mathrm{X}, \nu)\) 且 \(\varepsilon > 0\)。取 X 的一个紧子集 K，使得

\[
\int_ {\mathrm{X-K}} | \psi | d \nu <   \varepsilon .
\]

于是对 B 中任意 \(\varphi_{1}\) 和 \(\varphi_{2}\)，有

\[
\begin{array}{r l} & {\langle \varphi_ {1} - \varphi_ {2}, \psi \rangle = \int_ {\mathrm{X}} \psi (\varphi_ {1} - \varphi_ {2}) d \nu} \\ & {\qquad = \int_ {\mathrm{K}} \psi (\varphi_ {1} - \varphi_ {2}) d \nu + \int_ {\mathrm{X-K}} \psi (\varphi_ {1} - \varphi_ {2}) d \nu ,} \end{array}
\]

从而

\[
| \langle \varphi_ {1} - \varphi_ {2}, \psi \rangle | \leqslant \sup _ {x \in \mathrm{K}} | \varphi_ {1} (x) - \varphi_ {2} (x) | \| \psi \| _ {1} + 2 \mathrm{M} \varepsilon ,
\]

引理得证。

引理 6. --- 设 \(\psi \in \mathrm{L}^{1}(\mathrm{G})\)。令 B 是巴拿赫空间 \(\mathrm{L}^{\infty}(\mathrm{G})\) 的有界子集。当 B 赋予弱拓扑而映射 \(\varphi \mapsto \psi * \varphi\) 的值域 \(\mathcal{C}_b(\mathrm{G})\) 赋予紧收敛拓扑时，该映射从 B 到 \(\mathcal{C}_b(\mathrm{G})\) 连续。

设 \(\varphi\in\mathrm{L}^{\infty}(\mathrm{G})\) 且 \(\eta\in\mathrm{L}^{1}(\mathrm{G})\)。函数 \(\Delta^{-1}\check{\eta}\) 属于 \(\mathrm{L}^{1}(\mathrm{G})\)，且 \(\langle\eta,\check{\psi}\rangle=\langle\Delta^{-1}\check{\eta},\psi\rangle\)（参见 INT, VII, p. 19, § 1, n°3, 公式 (22)）。因此，映射 \(\varphi\mapsto\check{\varphi}\) 是赋予弱拓扑的空间 \(\mathrm{L}^{\infty}(\mathrm{G})\) 上的自同构。

设 \(\varphi\in\mathrm{L}^{\infty}(\mathrm{G})\)。由 INT, VIII, p. 167, § 4, no. 5, 命题 14，函数 \(\psi*\varphi\) 属于 \(\mathcal{C}_{b}(\mathrm{G})\)，且满足

\[
\begin{array}{r l} & {(\psi * \varphi) (g) = \int_ {\mathrm{G}} \psi (h) \varphi (h ^ {- 1} g) d \mu (h)} \\ & {\qquad = \int_ {\mathrm{G}} \psi (g y) \check {\varphi} (y) d \mu (y) = \langle \check {\varphi}, \pmb {\gamma} _ {\mathrm{G}} ^ {(1)} (g ^ {- 1}) \psi \rangle} \end{array}
\]

对所有 \(g \in G\) 都成立。因此，线性映射 \(u: \varphi \mapsto \psi * \varphi\) 从赋予弱拓扑的空间 \(L^{\infty}(G)\) 到赋予逐点收敛拓扑的 \(\mathcal{C}_{b}(G)\) 是连续映射。

设 \(\varphi\in\mathrm{B}\) 且 \((g,h)\in\mathrm{G}\times\mathrm{G}\)。由上式，有

\[
| u (\varphi) (g) - u (\varphi) (h) | \leqslant \| \check {\varphi} \| _ {\infty} \| (\boldsymbol {\gamma} _ {\mathrm{G}} ^ {(1)} (g ^ {- 1}) - \boldsymbol {\gamma} _ {\mathrm{G}} ^ {(1)} (h ^ {- 1})) \psi \| _ {1}.
\]

由于 B 有界且 G 在 \(L^{1}(G)\) 上的左正则表示连续（V, p. 405, no. 4），这意味着 \(u(B)\) 是 \(\mathcal{C}_{b}(G)\) 的等度连续子集。由前述结果及 TG, X, p. 16, 定理 1，断言得证。

引理 7. --- 设 \(\psi \in \mathrm{L}^1 (\mathbf{G})\) 满足 \(\psi \geqslant 0\) 且 \(\int \psi = 1\)。令 p 表示 \(\mathcal{C}_b(\mathbf{G})\) 上由 \(p(\varphi) = |\langle \varphi ,\psi \rangle |\) 对所有 \(\varphi \in \mathcal{C}_b(\mathbf{G})\) 定义的半范数。对每个 \(\varphi \in \mathrm{Pos}_1(\mathbf{G})\)，有

\[
\| \psi * \varphi - \varphi \| _ {\infty} \leqslant \sqrt {2 p (1 - \varphi)}.
\]

设 \(\varphi\in\mathrm{Pos}_{1}(\mathrm{G})\) 且 \(x\in\mathrm{G}\)。根据 INT, VIII, p. 167, § 4, no. 5, 命题 14，得到

\[
\begin{array}{r l} | \psi * \varphi (x) - \varphi (x) | = & \left| \int_ {\mathrm{G}} (\varphi (y ^ {- 1} x) - \varphi (x)) \psi (y) d \mu (y) \right| \\ & \leqslant \int_ {\mathrm{G}} | \varphi (y ^ {- 1} x) - \varphi (x) | \psi (y) d \mu (y) \\ & \leqslant \int_ {\mathrm{G}} \sqrt {2 (1 - \mathcal {R}   \varphi (y))} \psi (y) d \mu (y) \end{array}
\]

应用 V, p. 446 的界 (1)。于是由柯西–施瓦茨不等式可得

\[
\begin{array}{c} \| \psi * \varphi - \varphi \| _ {\infty} \leqslant \sqrt {2} \Bigl (\int_ {\mathrm{G}} (1 - \mathcal {R} (\varphi)) \psi   d \mu \Bigr) ^ {1 / 2} \Bigl (\int_ {\mathrm{G}} \psi   d \mu \Bigr) ^ {1 / 2} \\ \leqslant \sqrt {2} \sqrt {p (1 - \varphi)}, \end{array}
\]

从而得到结论。

引理 8. --- 设 K 是拓扑域，E 和 F 是 K 上的拓扑向量空间。设 f 是从 E 到 F 的映射，X 是 E 的子集，且 \(x \in X\)。若对于 F 中 0 的每个邻域 W，都存在 E 中 x 的一个邻域 U 以及从 X 到 F 的连续映射 g，使得映射 f \textbar{} X 在 x 处连续，当 \((f - g)(\mathrm{U} \cap \mathrm{X}) \subset \mathrm{W}\)。

设 \(W_{0}\) 是 F 中 0 的一个邻域，\(V_{0}\) 是 F 中 0 的一个对称邻域，满足 \(V_{0} + V_{0} + V_{0} \subset W_{0}\)。设 \(U_{0}\) 是 E 中 x 的一个邻域，g 是从 X 到 F 的连续映射，满足 \((f - g)(\mathrm{U}_{0} \cap \mathrm{X}) \subset \mathrm{V}_{0}\)。

存在 E 中 x 的一个邻域 \(U_{1}\)，使得 \(g(\mathrm{U}_{1} \cap \mathrm{X}) \subset g(x) + \mathrm{V}_{0}\)。对每个 \(y \in U_{0} \cap U_{1} \cap X\)，有

\[
\begin{array}{c} f (y) - f (x) = (f (y) - g (y)) + (g (y) - g (x)) + (g (x) - f (x)) \in V_{0} + V_{0} + V_{0} \subset W_{0}, \end{array}
\]

所以 \(f(y) \in f(x) + W_{0}\)，从而得到结论。

现在证明定理 3。以 \(\mathrm{Pos}_{1}(\mathrm{G})_{f}\)（分别以 \(\mathrm{Pos}_{1}(\mathrm{G})_{c}\)）表示赋予弱拓扑（分别赋予紧收敛拓扑）的集合 \(\mathrm{Pos}_{1}(\mathrm{G})\)；同样，以 \(\mathcal{C}_{b}(\mathrm{G})_{f}\)（分别以 \(\mathcal{C}_{b}(\mathrm{G})_{c}\)）表示赋予弱拓扑（分别赋予紧收敛拓扑）的空间 \(\mathcal{C}_{b}(\mathrm{G})\)。

从 \(\mathrm{Pos}_{1}(\mathrm{G})_{c}\) 到 \(\mathrm{Pos}_{1}(\mathrm{G})_{f}\) 的恒等映射连续（引理 5）。反之，以 \(\iota\) 表示 \(\mathrm{Pos}_{1}(\mathrm{G})_{f}\) 到 \(\mathcal{C}_{b}(\mathrm{G})_{c}\) 的包含映射。证明 \(\iota\) 连续即可完成证明。

设 \(\varphi_{1} \in \mathrm{Pos}_{1}(\mathbf{G})\)。为验证 \(\iota\) 在 \(\varphi_{1}\) 处连续，应用引理 8。设 W 是 \(\mathcal{C}_{b}(\mathbf{G})_{c}\) 中 0 的一个邻域。只需找到线性映射 \(u\)，从 \(\mathcal{C}_{b}(\mathbf{G})\) 到 \(\mathcal{C}_{b}(\mathbf{G})\)，它通过对子空间的传递诱导出从 \(\mathrm{Pos}_{1}(\mathbf{G})_{f}\) 到 \(\mathcal{C}_{b}(\mathbf{G})_{c}\) 的连续映射，并找到 \(\mathrm{Pos}_{1}(\mathbf{G})_{f}\) 中 0 的一个邻域 U，使得 \((\iota - u)(\mathrm{U} \cap \mathrm{Pos}_{1}(\mathbf{G})) \subset \mathrm{W}\)。

存在 \(\varepsilon > 0\)，使得 W 包含满足 \(\varphi \in \mathcal{C}_b(\mathrm{G})\) 且 \(\| \varphi \|_{\infty} \leqslant \varepsilon\) 的函数所成的集合。由于 \(\varphi_1(e) = 1\)，存在 G 中 \(e\) 的一个紧邻域 V，使得

\[
\sup _ {x \in \mathrm{V}} | 1 - \varphi_ {1} (x) | \leqslant \frac {\varepsilon^ {2}}{4}.
\]

令 \(\varphi_{\mathrm{V}}\) 为 V 的特征函数，并置 \(\psi_{\mathrm{V}} = \mu (\mathrm{V})^{-1}\varphi_{\mathrm{V}}\)。线性映射 \(u\colon \varphi \mapsto \psi_{\mathrm{V}}*\varphi\) 从 \(\mathrm{Pos}_1(\mathrm{G})_f\) 到 \(\mathcal{C}_b(\mathrm{G})_c\) 连续（引理 6）。

以 \(q_{V}\) 表示由 \(\varphi\mapsto|\langle\varphi,\psi_{V}\rangle|\) 在 \(\mathcal{C}_{b}(G)\) 上定义的半范数；它对于弱拓扑连续。由于 \(\psi_{V}\) 在 V 外为零，可得 \(q_{\mathrm{V}}(1-\varphi_{1})\leqslant\varepsilon^{2}/4\)；因此在 \(\varphi_{1}\) 的 \(\mathcal{C}_{b}(\mathrm{G})_{f}\) 中存在一个邻域 U，使得 \(q_{\mathrm{V}}(1-\varphi)\leqslant\varepsilon^{2}/2\) 对所有 \(\varphi\in U\) 成立。

设 \(\varphi\in\mathrm{U}\cap\mathrm{Pos}_{1}(\mathrm{G})\)。由引理 7，有

\[
\| (\iota - u) (\varphi) \| _ {\infty} \leqslant \sqrt {2 q _ {\mathrm{V}} (1 - \varphi)} \leqslant \varepsilon
\]

所以 \((\iota - u)(\mathrm{U}) \subset \mathrm{W}\)。定理得证。

一般而言，\(\operatorname{Pos}_{\leqslant1}(\mathbf{G})\) 上的弱拓扑不与紧收敛拓扑重合，群 R 的例子说明了这一点。

回忆我们以 \(\widehat{G}\) 表示 \(G\) 的不可约酉表示的等价类集合（V, p. 393, 定义 8）。

定理 4（盖尔范德--赖科夫）。--- 对 G 中每个 \(x \neq e\)，存在 G 的不可约酉表示 \(\pi\)，使得 \(\pi(x)\) 不是 \(\pi\) 的空间上的恒等自同态。

设 \(x \in G\) 满足 \(\pi(x)\) 对每个 \(\pi \in \widehat{G}\) 都是恒等映射。则有 \(\varphi(x) = \varphi(e) = 1\) 对每个极点 \(\varphi\)（属于 \(\mathrm{Pos}_{1}(\mathrm{G})\)）（V, p. 444, 命题 11）；又由于命题 14，这对每个 \(\varphi \in \mathrm{Pos}_{1}(\mathrm{G})\) 都成立，因为在线性泛函 \(\varphi \mapsto \varphi(e)\) 在赋予弱拓扑的 \(\mathrm{Pos}_{1}(\mathrm{G})\) 上是连续的（定理 3）。但是若 \(x \neq e\)，则存在范数为 1 的 \(f \in L^{2}(\mathrm{G})\)，使得 \(\langle f | \boldsymbol{\gamma}_{\mathrm{G}}(x)f \rangle = 0\)（V, p. 406, 引理 3）。由于这个等式左端具有 \(\varphi(x)\) 的形式，其中 \(\varphi \in \mathrm{Pos}_{1}(\mathrm{G})\)，故矛盾。

若 G 不是局部紧的，则 G 的不可约酉表示未必能分离 G 的点。

推论。--- 对 G 中任意 \(g \neq h\)，存在不可约表示 \(\pi \in \widehat{G}\) 及 \(\pi\) 的一个矩阵系数 f，使得 \(f(g) \neq f(h)\)。特别地，\(\Upsilon(\mathrm{G})\) 是 \(\mathcal{C}_{b}(\mathrm{G})\) 的能分离 G 的点的子代数。

设 \(g \neq h\) 属于 G。存在不可约表示 \(\pi \in \widehat{G}\)，使得 \(\pi(g) \neq \pi(h)\)（定理 4），所以在 \(\pi\) 的空间中存在 x 和 y，使得 \(\langle x | \pi(g)y \rangle \neq \langle x | \pi(h)y \rangle\)。

\subsection{局部紧交换群上的正定函数}

在本节中，G 是局部紧交换群，\(\mu\) 表示 G 上的哈尔测度。以 \(\widehat{G}\) 表示 G 的对偶群，以 \(\widehat{\mu}\) 表示 \(\mu\) 的对偶哈尔测度（II, p. 214, 定义 4）。以 F 表示 G 上有界测度的巴拿赫空间 \(\mathcal{M}^{1}(G)\) 上的傅里叶变换。

由于 \(\mathcal{C}_0(\mathrm{G})\) 是 \(\mathcal{K}(\mathrm{G})\) 在 \(\mathcal{C}_b(\mathrm{G})\) 中的闭包，巴拿赫空间 \(\mathcal{M}^1 (\mathrm{G})\)（\(\mathcal{K}(\mathrm{G})\) 赋予 \(\mathcal{C}_b(\mathrm{G})\) 的拓扑的对偶）与 \(\mathcal{C}_0(\mathrm{G})\) 的对偶相等同（参见 INT, III, p. 56, § 1, no. 8）。在 \(\mathcal{M}^1 (\mathrm{G})\) 上赋予与此对偶性相联系的弱拓扑。

引理 9. --- 傅里叶变换是从赋予弱拓扑的 \(\mathcal{M}^{1}(\mathrm{G})\) 到 \(\mathcal{C}_{b}(\widehat{\mathrm{G}})\) 的连续映射，后者赋予由弱拓扑 \(\sigma(\mathrm{L}^{\infty}(\widehat{\mathrm{G}}),\mathrm{L}^{1}(\widehat{\mathrm{G}}))\) 诱导的拓扑。

设 \(\mathfrak{F}\) 是 \(\mathcal{M}^1 (\mathrm{G})\) 上弱收敛于 G 上有界测度 \(\nu\) 的滤子。对每个 \(\varphi \in \mathrm{L}^1 (\widehat{\mathrm{G}})\)，有 \(\mathcal{F}(\varphi)\in \mathcal{C}_0(\mathrm{G})\)（II, p. 209, 命题 4），从而

\[
\begin{array}{r l} & {\int_ {\widehat {\mathrm{G}}} \varphi \mathcal {F} (\nu) d \widehat {\mu} = \int_ {\mathrm{G}} \mathcal {F} (\varphi) d \nu} \\ & {\qquad = \lim _ {\eta , \mathfrak {F}} \int_ {\mathrm{G}} \mathcal {F} (\varphi) d \eta = \lim _ {\eta , \mathfrak {F}} \int_ {\widehat {\mathrm{G}}} \varphi \mathcal {F} (\eta) d \widehat {\mu},} \end{array}
\]

由傅里叶变换的转置性质（II, p. 221, 命题 13）。引理得证。

定理 5（博赫纳）。--- 连续函数 \(\varphi\) 在 \(\widehat{\mathbf{G}}\) 上且属于 \(\operatorname{Pos}(\widehat{\mathbf{G}})\)，当且仅当存在 \(\nu\) 在 \(\mathbf{G}\) 上的有界正测度，使得 \(\varphi = \mathcal{F}(\nu)\)。

设 \(\nu \in \mathcal{M}^{1}(\mathrm{G})\) 是正测度。其傅里叶变换连续。对每个有限族 \((x_{i})_{i\in \mathrm{I}}\)（属于 \(\widehat{\mathbf{G}}\)）以及复数的每个有限族 \((t_i)_{i\in \mathrm{I}}\)，可得

\[
\begin{array}{r} \sum_ {i \in \mathrm{I}} \sum_ {j \in \mathrm{I}} \bar {t} _ {i} t _ {j} \mathcal {F} (\nu) (x _ {i} ^ {- 1} x _ {j}) = \sum_ {i \in \mathrm{I}} \sum_ {j \in \mathrm{I}} \bar {t} _ {i} t _ {j} \int_ {\mathrm{G}} x _ {i} \overline {{x}} _ {j} d \nu \\ = \int_ {\mathrm{G}} \left| \sum_ {i \in \mathrm{I}} \bar {t} _ {i} x _ {i} \right| ^ {2} d \nu \geqslant 0, \end{array}
\]

由于 \(\nu\) 是正测度。因此，\(\nu\) 的傅里叶变换是 \(\hat{G}\) 上的正定函数（V, p. 432, 定理 1, (ii)）。

证明逆向蕴含。给集合 \(\operatorname{Pos}_{\leqslant1}(\widehat{\mathbf{G}})\) 赋予弱拓扑，该拓扑像前面一样由弱拓扑 \(\sigma(\mathrm{L}^{\infty}(\widehat{\mathrm{G}}),\mathrm{L}^{1}(\widehat{\mathrm{G}}))\) 诱导。集合 \(\operatorname{Pos}_{\leqslant1}(\widehat{\mathbf{G}})\) 是紧凸集（V, p. 448, 命题 13 的推论）。由 V, p. 450, 命题 14，V, p. 444, 命题 11 以及 V, p. 390, 推论 7，\(\operatorname{Pos}_{\leqslant1}(\widehat{\mathbf{G}})\) 的极点是零函数和 \(\widehat{G}\) 中的元素。

令 N 为 G 上质量 \(\leqslant 1\) 的正有界测度的集合；它在赋予弱拓扑的 \(\mathcal{M}^{1}(G)\) 中是紧的（EVT, III, p. 17, 推论 3）。由引理 9 及证明的第一部分，G 上的傅里叶变换通过对子空间的传递定义了从 \(\mathcal{N}\) 到 \(\mathrm{Pos}_{\leqslant 1}(\widehat{\mathrm{G}})\) 的连续映射；由齐次性，只需证明此映射是满射。\(\mathcal{F}(\mathcal{N})\) 是傅里叶变换下 \(\mathcal{N}\) 的像，是紧凸集；它包含零函数和 \(\widehat{\mathrm{G}}\) 中的元素（事实上，由 II, p. 220, 定理 2，这些元素都具有 \(\mathrm{ev}_x\colon \chi \mapsto \chi (x)\) 的形式，其中 \(x\) 是 G 中某个元，且 \(\mathrm{ev}_x = \mathcal{F}(\varepsilon_{x^{-1}})\)）。因此集合 \(\mathcal{F}(\mathcal{N})\) 包含 \(\mathrm{Pos}_{\leqslant 1}(\mathrm{G})\) 的极点，由克赖因–米尔曼定理（EVT, II, p. 59, 定理 1）可得 \(\mathcal{F}(\mathcal{N}) = \mathrm{Pos}_{\leqslant 1}(\mathrm{G})\)。证明完毕。

当 \(G = R^{k}\)，其中整数 \(k \in N\) 时，本定理对应于 INT, IX, p. 94, § 6, no. 12, 命题 11。

\section{紧群的表示}

在本段中，除非另有明确说明，所考虑的所有拓扑向量空间都是复的。

固定一个紧拓扑群 G，其单位元记为 e。群 G 是单模的（INT, VII, p. 20, § 1, n° 3, 命题 3 的推论）。以 μ 表示 G 上的归一化哈尔测度（即满足 μ(G) = 1 的测度），并且对 \(1 \leqslant p \leqslant +\infty\)，以 \(\mathcal{L}^{p}(G)\)（分别以 \(\mathrm{L}^{p}(\mathrm{G})\)）表示空间 \(\mathcal{L}_{\mathbf{C}}^{p}(\mathrm{G}, \mu)\)（分别表示空间 \(\mathrm{L}_{\mathbf{C}}^{p}(\mathrm{G}, \mu)\)）。卷积总是相对于测度 μ 取。设 \(p \in [1, +\infty]\)。将 \(\mathcal{C}(\mathrm{G})\) 与 \(\mathrm{L}^{p}(\mathrm{G})\) 的一个子空间等同，这是允许的，因为 μ 的支集等于 G。

对每个不可约酉表示 \(\pi \in \widehat{\mathbf{G}}\)，令 \(\mathrm{E}_{\pi}\) 为 \(\pi\) 的空间。

\subsection{有限维表示的半单性}

回忆（INT, VII, p. 71, § 3, n° 1, 引理 1）：对于 G 在希尔伯特空间 E 上的每个连续表示 \(\varrho\)，E 上存在非退化正埃尔米特形式 \(q\)，使得由 q 定义的 E 的拓扑向量空间结构与 E 的原结构相同，并且 \(\varrho\) 是 G 在赋予内积 q 的希尔伯特空间 E 上的酉表示。

命题 1. --- 设 \(\varrho\) 是 G 在分离拓扑向量空间 E 上的有限维线性表示。E 上存在内积 \(q\)，使得 \(\varrho\) 是赋予 \(q\) 的希尔伯特空间 E 上的酉表示。特别地，\(\varrho\) 是半单的。

由于 E 是有限维的，E 上存在希尔伯特空间结构。应用前述注记以及 V, p. 392, 推论 2 即得结果。

注。--- 稍后我们将看到（V, p. 466, 推论 4），G 的每个酉表示都是半单的。

推论。--- 设 \(\varrho_{1}\) 和 \(\varrho_{2}\) 是 G 的有限维表示。表示 \(\varrho_{1}\) 和 \(\varrho_{2}\) 同构，当且仅当它们的特征标相等。

这由命题及 V, p. 392, 推论 3 得出。

\subsection{不可约表示}

引理 1. --- G 的每个不可约酉表示都是平方可积的。

事实上，不可约酉表示的矩阵系数连续且有界，所以在 G 上平方可积，因为 G 是紧的。

命题 2. --- 设 π 是 G 在希尔伯特空间 E 上的不可约酉表示。E 的维数有限，且等于 π 的形式次数。

表示 \(\pi\) 是平方可积的（引理 1），所以形式次数 \(c_{\mu}(\pi)\) 相对于 \(\mu\) 有定义（V, p. 423, 定义 4）；它是严格为正的实数。设 \((e_i)_{i \in I}\) 是 E 中的有限标准正交族。设 \(x\) 是 E 中范数为 1 的元（这样的元存在，因为 E 非零）。对 \(i \in I\)，有

\[
c _ {\mu} (\pi) \int_ {\mathrm{G}} | \langle e _ {i} | \pi (g) x \rangle | ^ {2} d \mu (g) = 1
\]

（V, p. 424, 命题 8）。对 \(i \in I\) 将此公式求和，得到

\[
\begin{array}{r l} & {\mathrm{Card(I)} = c _ {\mu} (\pi) \int_ {\mathrm{G}} \sum_ {i \in \mathrm{I}} | \langle e _ {i} | \pi (g) x \rangle | ^ {2} d \mu (g)} \\ & {\qquad \leqslant c _ {\mu} (\pi) \int_ {\mathrm{G}} \| \pi (g) x \| ^ {2} d \mu (g) = c _ {\mu} (\pi)} \end{array}
\]

由贝塞尔不等式（EVT, V, p. 21, 命题 4）。因此 I 的基数上界为 \(c_{\mu}(\pi)\)。这意味着 E 的维数有限。

于是可将前述结果应用于 E 的标准正交基 \((e_{i})_{i\in\mathrm{I}}\)。由 EVT, V, p. 22, 命题 5，得到

\[
\begin{array}{r} \dim (\mathrm{E}) = c _ {\mu} (\pi) \int_ {\mathrm{G}} \sum_ {i \in \mathrm{I}} | \langle e _ {i} | \pi (g) x \rangle | ^ {2} d \mu (g) \\ = c _ {\mu} (\pi) \int_ {\mathrm{G}} \| \pi (g) x \| ^ {2} d \mu (g) = c _ {\mu} (\pi), \end{array}
\]

证明完毕。

特别地，因此可以谈论特征标 \(\chi_{\pi}\)，其中 \(\pi\) 是 G 的一个不可约酉表示。它是 G 上的连续函数。

推论 1. --- 设 \(\pi\) 是 G 在希尔伯特空间 E 上的不可约酉表示。\(\pi\) 的特征标是对合代数 \(\mathrm{L}^{1}(\mathrm{G})\) 的埃尔米特元。

设 \(x \in G\)。由于 \(G\) 是单模的，有 \(\chi_{\pi}^{*}(x) = \overline{\operatorname{Tr}(\pi(x^{-1}))}\)，又有 \(\chi_{\pi}^{*}(x) = \operatorname{Tr}(\pi(x))\)，这是由于 \(\pi\) 是酉表示。

推论 2. --- a) 设 \(\pi\) 是 G 在希尔伯特空间 E 上的不可约酉表示。有

\[
\int_ {\mathrm{G}} \overline {{\langle x \mid \pi (g) x ^ {\prime} \rangle}} \langle y \mid \pi (g) y ^ {\prime} \rangle d \mu (g) = \frac {1}{\dim (\mathrm{E})} \overline {{\langle x \mid y \rangle}} \langle x ^ {\prime} \mid y ^ {\prime} \rangle
\]

对所有 \((x,y,x',y')\in \mathrm{E}^4\) 都成立。

\begin{enumerate}
\def\labelenumi{\alph{enumi})}
\setcounter{enumi}{1}
\tightlist
\item
  设 \(\pi_{1}\)（分别为 \(\pi_{2}\)）是 G 在希尔伯特空间 \(E_{1}\)（分别为 \(E_{2}\)）上的不可约酉表示。若 \(\pi_{1}\) 和 \(\pi_{2}\) 不同构，则有
\end{enumerate}

\[
\int_ {G} \overline {{\langle x \mid \pi_ {1} (g) x ^ {\prime} \rangle}} \langle y \mid \pi_ {2} (g) y ^ {\prime} \rangle d \mu (g) = 0
\]

对所有 \((x,x^{\prime},y,y^{\prime})\in \mathrm{E}_{1}^{2}\times \mathrm{E}_{2}^{2}\) 都成立

这立即由 V, p. 424, 命题 8 和 V, p. 424, 命题 9 得出，同时考虑引理 1 及公式 \(c_{\mu}(\pi) = \dim(\mathrm{E})\)（命题 2）。

推论 3. --- 设 \(\pi_1\) 和 \(\pi_2\) 是 G 的不可约表示。在希尔伯特空间 \(\mathrm{L}^{2}(\mathrm{G})\) 中，有 \(\langle \chi_{\pi_1} | \chi_{\pi_2} \rangle = 1\)（当 \(\pi_1\) 和 \(\pi_2\) 同构时），否则有 \(\langle \chi_{\pi_1} | \chi_{\pi_2} \rangle = 0\)。换言之，类 \(\pi \in \widehat{\mathbf{G}}\) 的特征标所成的族是 \(\mathrm{L}^{2}(\mathrm{G})\) 中的标准正交族。

这由 V, p. 457, 命题 1 及推论 2 得出；注意，对于 G 的有限维酉表示的空间的每个标准正交基 \((e_{i})_{i\in\mathrm{I}}\)，以及表示 \(\pi\) 和每个 \(g\in G\)，都有公式

\[
\chi_ {\pi} (g) = \sum_ {i \in I} \langle e _ {i} | \pi (g) e _ {i} \rangle .
\]

推论 4. --- a) 设 \(\varrho\) 是 G 的连续有限维表示。有

\[
\chi_ {\varrho} = \sum_ {\pi \in \widehat {G}} \dim (\operatorname{Hom} _ {G} (\pi , \varrho)) \chi_ {\pi} = \sum_ {\pi \in \widehat {G}} \dim (\operatorname{Hom} _ {G} (\varrho , \pi)) \chi_ {\pi}.
\]

\begin{enumerate}
\def\labelenumi{\alph{enumi})}
\setcounter{enumi}{1}
\tightlist
\item
  设 \(\varrho_{1}\) 和 \(\varrho_{2}\) 是 G 的连续有限维表示。有
\end{enumerate}

\[
\left\langle \chi_ {\varrho_ {1}} \mid \chi_ {\varrho_ {2}} \right\rangle = \dim (\operatorname{Hom} _ {\mathrm{G}} (\varrho_ {1}, \varrho_ {2})) = \dim (\operatorname{Hom} _ {\mathrm{G}} (\varrho_ {2}, \varrho_ {1})).
\]

\begin{enumerate}
\def\labelenumi{\alph{enumi})}
\setcounter{enumi}{2}
\tightlist
\item
  G 的连续有限维表示 \(\varrho\) 不可约，当且仅当 \(\| \chi_{\varrho}\|^{2} = 1\)。
\end{enumerate}

由于 \(\varrho\) 是半单的，它同构于其 \(\pi\)-同型分量 \(\mathrm{M}_{\pi}(\varrho)\) 的直和，其中 \(\pi \in \widehat{\mathrm{G}}\)。以 \(m_{\pi}(\varrho)\) 表示 \(\pi\) 在 \(\varrho\) 中的重数（V, p. 398, 定义 10）；于是有

\[
\chi_ {\varrho} = \sum_ {\pi \in \widehat {G}} m _ {\pi} (\varrho) \chi_ {\pi}.
\]

因此，断言 a) 由以下事实得出：

\[
m _ {\pi} (\varrho) = \dim \operatorname{Hom} _ {\mathrm{G}} (\pi , \varrho) = \dim \operatorname{Hom} _ {\mathrm{G}} (\varrho , \pi)
\]

（V, p. 377, 公式 (1) 及 V, p. 387, 推论 2）。

由特征标的双线性和标准正交性，断言 \(a\) ) 蕴含

\[
\langle \chi_ {\varrho_ {1}} \mid \chi_ {\varrho_ {2}} \rangle = \sum_ {\pi \in \widehat {\mathrm{G}}} \dim (\operatorname{Hom} _ {\mathrm{G}} (\pi , \varrho_ {1})) \dim (\operatorname{Hom} _ {\mathrm{G}} (\pi , \varrho_ {2})),
\]

另一方面，有典范同构

\[
\begin{array}{c} \operatorname{Hom} _ {\mathrm{G}} (\varrho _ {1}, \varrho _ {2}) \to \bigoplus_ {(\pi _ {1}, \pi _ {2}) \in \widehat {\mathrm{G}} \times \widehat {\mathrm{G}}} \operatorname{Hom} _ {\mathrm{G}} (\mathrm{M} _ {\pi _ {1}} (\varrho _ {1}), \mathrm{M} _ {\pi _ {2}} (\varrho _ {2})) \\ \qquad \qquad \qquad \qquad \qquad \qquad \qquad \qquad \qquad \qquad \qquad \qquad \qquad \qquad \qquad \qquad \qquad \qquad \qquad \qquad \qquad \qquad \qquad \qquad \qquad \qquad \qquad \qquad \qquad \qquad \qquad \qquad \qquad \qquad \qquad \to \bigoplus_ {\pi \in \widehat {\mathrm{G}}} \operatorname{Hom} _ {\mathrm{G}} (\mathrm{M} _ {\pi} (\varrho _ {1}), \mathrm{M} _ {\pi} (\varrho _ {2})) \end{array}
\]

（V, p. 377, 公式 (1)），从而可得

\[
\dim (\operatorname{Hom} _ {\mathrm{G}} (\varrho_ {1}, \varrho_ {2})) = \sum_ {\pi \in \widehat {\mathrm{G}}} m _ {\pi} (\varrho_ {1}) m _ {\pi} (\varrho_ {2}),
\]

从而得到断言 \(b\)。断言 \(c\) 也由 \(a\) 和舒尔引理（V, p. 386, 命题 6）得出。

推论 5. --- 设 \(\pi_1\) 和 \(\pi_2\) 是 G 的不可约酉表示。则 \(\pi_1(\overline{\chi}_{\pi_2}) = 0\)；若 \(\pi_1\) 不同构于 \(\pi_2\)，而 \(\pi_1(\overline{\chi}_{\pi_1})\) 是乘以 \(1/\dim(\pi_1)\) 的映射。

\(\pi_2\) 的特征标是 G 上的连续中心函数，因此线性映射 \(u = \pi_1(\overline{\chi}_{\pi_2})\) 有定义且属于空间 \(\mathrm{Hom}_{\mathrm{G}}(\pi_1,\pi_1)\)。由舒尔引理（V, p. 386, 命题 6），它是一个倍乘映射，其迹为

\[
\begin{array}{c} \operatorname{Tr} (u) = \operatorname{Tr} \Bigl (\int_ {\mathrm{G}} \overline {{\chi_ {\pi_ {2}} (g)}}   \pi_ {1} (g) d \mu (g) \Bigr) \\ = \int_ {\mathrm{G}} \overline {{\chi_ {\pi_ {2}} (g)}} \chi_ {\pi_ {1}} (g) d \mu (g). \end{array}
\]

由正交关系，\(u\) 的迹为零；若 \(\pi_1\) 不同构于 \(\pi_2\)，否则等于 1。断言得证。

注。--- 当 G 有限时，舒尔正交关系（分别为特征标正交公式）与 A, VIII, p. 399 中的公式（分别与 A, VIII, p. 400, 命题 4 中的公式）一致。另一方面，有限群特征标的“第二正交关系”（A, VIII, p. 402, 公式 (32)）在 G 紧时没有完全对应的形式。

当 G 有限时，与 e 的共轭类对应的第二正交公式的特殊情形是公式

\[
\frac {1}{\mathrm{CardG}} \sum_ {\pi \in \widehat {\mathrm{G}}} \dim (\pi) \chi _ {\pi} (g) = \left\{ \begin{array}{l l} 1 & \text {if} g = e \\ 0 & \text {otherwise}, \end{array} \right.
\]

这也可解释为计算 G 的正则表示 \(\gamma_{G}\) 的特征标，并且等价于对从 G 到 C 的每个函数 f 都有公式 \(\mathrm{Tr}(\boldsymbol{\gamma}_{\mathrm{G}}(f)) = f(e)\)。

再设 G 是任意紧群。对每个函数 \(f \in \mathcal{C}(\mathrm{G})\)，\(\boldsymbol{\gamma}_{\mathrm{G}}(f)\) 是 \(\mathrm{L}^{2}(\mathrm{G})\) 的自同态，与自同态 \(\varphi \mapsto f * \varphi\) 重合，且具有有限迹（V, p. 407, 引理 4 及 III, p. 33, 推论 2）。由同上及 IV, p. 177, 定理 2，还有

\[
\operatorname{Tr} \left(\boldsymbol {\gamma} _ {\mathrm{G}} (f)\right) = \int_ {\mathrm{G}} f \left(x ^ {- 1} x\right) d \mu (x) = f (e).
\]

命题 3. --- 设 \(G = G_{1} \times G_{2}\)，其中 \(G_{1}\) 和 \(G_{2}\) 是紧群。从 \(\widehat{G}_{1} \times \widehat{G}_{2}\) 到 \(\widehat{G}\) 的映射 b 将 \((\pi_{1}, \pi_{2})\) 送到 \(\pi_{1} \boxtimes \pi_{2}\) 的类，且 b 是双射。

由 V, p. 389, 推论 6，映射 b 定义良好。

设 \(\pi_1\) 和 \(\pi_2\) 是 \(\widehat{\mathbf{G}}\) 的元素。设 \(\pi \in \widehat{\mathbf{G}}\)。\(\pi\)-同型分量是限制 \(\varpi\) 的 \(\pi_1 \boxtimes \pi_2\) 在 \(\mathrm{G}_1 \times \{e\}\) 上，其等于 \(\varpi\) 当 \(\pi_1 = \pi\) 时，否则为零（V, p. 384, 引理 8）。因此，酉表示 \(\pi_1 \boxtimes \pi_2\) 在同构意义下确定 \(\pi_1\)；同样，它确定 \(\pi_2\)，这证明 \(b\) 是单射。

证明 \(b\) 是满射。设 \(\pi\) 是 \(G_1 \times G_2\) 在希尔伯特空间 E 上的不可约酉表示。它是有限维的（命题 2）。设 \(\pi_1\) 是 \(G_1\) 在希尔伯特空间 \(E_1\) 上的不可约酉表示，使得 \(\pi\) 到 \(G_1 \times \{e\}\) 的限制包含一个与 \(\pi_1\) 同构的子表示（V, p. 379, 引理 3）。令 \(F = \text{Hom}_{G_1}(\pi_1, \pi)\)。这是非零有限维向量空间。对 \(h \in G_2\) 和 \(u \in F\)，令 \(\varrho(h)(u)\) 为从 \(E_1\) 到 E 的线性映射，由 \(x \mapsto \pi(e, h)(u(x))\) 定义。由于 \(G_1 \times \{e\}\) 与 \(\{e\} \times G_2\) 在 G 中交换，映射 \(\varrho(h)(u)\) 属于空间 F。映射 \(\varrho\) 是 \(G_2\) 在空间 F 上的线性表示；它是连续的。由于 F 非零，存在不可约子表示 \(\pi_2\) 的 \(\varrho\)（V, p. 379, 引理 3）。令 \(E_2\) 为 \(\pi_2\) 的空间。于是，线性映射 \(v: E_1 \otimes E_2 \to E\) 满足 \(x \otimes u \mapsto u(x)\) 对 \(x \in E_1\) 和 \(u \in E_2\) 成立，是从 \(\pi_1 \boxtimes \pi_2\) 到 \(\pi\) 的非零 G-态射；由于表示 \(\pi\) 和 \(\pi_1 \boxtimes \pi_2\) 都不可约且有限维，态射 \(v\) 是同构（V, p. 378, 引理 2）。

本命题可与 A, VIII, p. 208 的定理 1 比较。

\subsection{彼得--外尔定理}

回忆我们以 \(\Theta(\mathrm{G})\) 表示 G 的有限维酉表示的矩阵系数空间。空间 \(\Theta(\mathrm{G})\) 是 \(\mathcal{C}(\mathrm{G})\) 的一个在复共轭下稳定的含单位元子代数（V, p. 386, 命题 5）。

对 \(\pi \in \widehat{\mathrm{G}}\)，令 \(\varrho_{\mathrm{G}}(\pi)\) 为 \(\mathcal{C}(\mathrm{G})\) 中由 \(\pi\) 的矩阵系数生成的子空间。将其与 \(\mathrm{L}^2 (\mathrm{G})\) 的一个子空间等同。

空间 \(\Theta(G)\) 等于各空间 \(\varrho_{G}(\pi)\) 的和，其中 \(\pi\in\widehat{G}\)。事实上，\(\Theta(G)\) 的每个元都是 G 的不可约表示的矩阵系数之和，因为由 V, p. 457, 命题 1，G 的有限维表示是半单的。此外，该和是直和，因为由 V, p. 458, 推论 2，各空间 \(\varrho_{G}(\pi)\) 两两正交。

命题 4. — 空间 \(\Theta(\mathrm{G})\) 在 \(\mathcal{C}(\mathrm{G})\) 和 \(\mathrm{L}^2(\mathrm{G})\) 中都稠密。

含单位元子代数 \(\Theta(\mathrm{G})\) 的 \(\mathcal{C}(\mathrm{G})\) 在共轭下稳定。由 V, p. 457, 命题 2，它与子代数 \(\Upsilon(\mathrm{G})\) 重合，故能分离 G 的点（V, p. 454, 定理 4 的推论）。因此，由 TG, X, p. 40, 推论 2，它在 \(\mathcal{C}(\mathrm{G})\) 中稠密，从而在 \(\mathrm{L}^2(\mathrm{G})\) 中也稠密（参见 INT, IV, p. 155, §4, no. 7, 命题 13）。

回忆 G 的双正则酉表示是表示 \(\varrho_{G}\)，它从 G × G 到 \(\mathrm{L}^{2}(\mathrm{G})\)，满足 \((g_{1}, g_{2}) \mapsto \gamma_{\mathrm{G}}(g_{1}) \delta_{\mathrm{G}}(g_{2})\)。

命题 5. --- 设 \(\pi \in \widehat{\mathbf{G}}\)。空间 \(\varrho_{\mathrm{G}}(\pi)\) 是 \(\varrho_{\mathrm{G}}\) 的一个子表示，与 \(\overline{\pi} \boxtimes \pi\) 同构。特别地，它是 \(\mathrm{G} \times \mathrm{G}\) 的不可约表示。

由 V, p. 422, 命题 7（应用 \(Z = \{e\}\)），空间 \(\varrho_{G}(\pi)\) 是 \(\varrho_{G}\) 的子表示，且从 \(\overline{\mathrm{E}}_{\pi} \otimes \mathrm{E}_{\pi}\) 到 \(L^2(G)\)、将 \(x \otimes y\) 送到矩阵系数 \(g \mapsto \langle x | \pi(g)y \rangle\) 的线性映射，是一个 \((G \times G)\)-同构，从 \(\overline{\pi} \boxtimes \pi\) 到 \(\varrho_{G}(\pi)\)。因此子表示 \(\varrho_{G}(\pi)\) 是不可约的（V, p. 461, 命题 3）。

定理 1（彼得--外尔）。--- 设 G 是紧拓扑群。G 的双正则表示 \(\varrho_{\mathrm{G}}\) 是各子表示 \(\varrho_{\mathrm{G}}(\pi)\)（其中 \(\pi \in \widehat{\mathrm{G}}\)）的希尔伯特直和。

各空间 \(\varrho_{\mathrm{G}}(\pi)\) 两两正交；它们是 G 的双正则表示的不可约子表示（命题 5），并且两两不同构（V, p. 461, 命题 3）。

定理即由以下事实得出：空间 \(\Theta(\mathrm{G})\) 是各空间 \(\varrho_{\mathrm{G}}(\pi)\) 的和，其中 \(\pi\) 遍历 \(\widehat{\mathrm{G}}\)，且该和在 \(\mathrm{L}^2(\mathrm{G})\) 中稠密（命题 4）。

推论 1. --- 对每个 \(\pi\in\widehat{G}\)，子表示 \(\varrho_{\mathrm{G}}(\pi)\) 是 \((\overline{\pi}\boxtimes\pi)\)-同型分量的 \(\varrho_{\mathrm{G}}\)。

令 F 为 \((\overline{\pi} \boxtimes \pi)\)-同型分量的 \(\varrho_{G}\)。空间 F 包含 \(\varrho_{\mathrm{G}}(\pi)\)（命题 5）。此外，对每个 \(\tau \in \widehat{G}\)（不同于 \(\pi\)），F 与 \(\varrho_{\mathrm{G}}(\tau)\) 的交为零（V, p. 461, 命题 3）。应用定理可推出 \(\mathrm{F} = \varrho_{\mathrm{G}}(\pi)\)。

推论 2. --- 设 \(\pi_1\) 和 \(\pi_2\) 是不同构的 G 的不可约表示。\((\overline{\pi}_1 \boxtimes \pi_2)\)-同型分量的 \(\varrho_G\) 为零。

推论 3. --- G 的右（分别为左）正则表示同构于希尔伯特直和

\[
\bigoplus_ {\pi \in \widehat {\mathrm{G}}} \pi^ {\dim (\pi)}.
\]

由定理及命题 5，G 的双正则表示限制到子群 \(\mathrm{H} = \{e\} \times \mathrm{G}\) 后，同构于各表示 \(\overline{\pi} \boxtimes \pi\) 在 H 上的限制的希尔伯特直和，其中 \(\pi \in \widehat{\mathrm{G}}\)。这些表示同构于 \(\dim(\pi)\) 个 \(\pi\) 的直和（V, p. 384, 引理 8）。这给出右正则表示的结果，左正则表示的情形类似。

推论 4. --- 设 \(\pi\in\widehat{G}\)。\(\pi\)-同型分量的 \(\delta_{G}\)（分别为 \(\overline{\pi}\)-同型分量的 \(\gamma_{G}\)）等于 \(\varrho_{G}(\pi)\)。

论证与前一推论类似，只需将 \(\varrho_{\mathrm{G}}(\pi)\) 看作 \(\delta_{\mathrm{G}}\)（分别为 \(\gamma_{\mathrm{G}}\)）的子表示。

注。--- 1) 若 G 有限，这些陈述对应于 A, VIII, p. 398 的注记中的结果。

\begin{enumerate}
\def\labelenumi{\arabic{enumi})}
\setcounter{enumi}{1}
\item
  假设 G 是交换的。集合 \(\widehat{G}\) 与 G 的对偶群相等同（V, p. 393, 注记）。对每个 \(\chi \in \widehat{G}\)，空间 \(\varrho_{\mathrm{G}}(\chi)\) 是 \(\mathrm{L}^2 (\mathrm{G})\) 中由 \(\chi\) 生成的一维子空间。因此，G 的定理 1 等价于 II, p. 215 的定理 1 的推论。
\item
  若存在整数 \(n \geqslant 0\)，使得 G 是 \(\mathbf{GL}(\mathbf{C}^{n})\) 的紧子群，则 G 在 \(\mathbf{GL}(\mathbf{C}^{n})\) 中的恒等表示足以分离 G 的点，随后可以直接证明命题 4，再证明彼得–外尔定理，而无需援引盖尔范德–赖科夫定理。
\item
  彼得–外尔定理特别蕴含，从 G 到 \(\prod_{\pi\in\widehat{G}}\mathbf{U}(\mathrm{E}_{\pi})\) 的自然连续同态是单射。
\end{enumerate}

\subsection{矩阵系数与 G-有限函数}

命题 6. --- \(\mathrm{L}^{2}(\mathrm{G})\) 的下列子空间相等：

\begin{enumerate}
\def\labelenumi{\alph{enumi})}
\item
  空间 \(\Theta (\mathbf{G})\)
\item
  \(\varrho_{\mathrm{G}}(\pi)\) 的 \(\mathrm{L}^2 (\mathrm{G})\) 子空间的代数直和。
\item
  \(\gamma_{\mathrm{G}}\) 的 G-有限向量空间（参见 V, p. 376）；
\item
  \(\delta_{G}\) 的 G-有限向量空间；
\item
  \((\mathbf{G} \times \mathbf{G})\)-有限向量空间的 \(\varrho_{\mathrm{G}}\)。
\end{enumerate}

特别地，\(\gamma_{\mathrm{G}}\)、\(\delta_{G}\) 或 \(\varrho_{\mathrm{G}}\) 的每个 G-有限向量都属于 \(\mathcal{C}(\mathrm{G})\)。

以 \(F_{a}\)（分别以 \(F_{b}\)、\(F_{c}\)、\(F_{d}\)、\(F_{e}\)）表示由条件 a)（分别由 b)、c)、d)、e)）定义的空间。我们已经注意到 \(F_{a} = F_{b}\)。

有 \(F_{b} \subset F_{c}\)，因为 \(\varrho_{\mathrm{G}}(\pi)\) 是 \(\gamma_{G}\) 的有限维子表示，对所有 \(\pi \in \widehat{G}\) 成立。反之，设 f 是 \(\gamma_{G}\) 的 G-有限向量。子空间 \(E_{f}\) 由函数 \(\gamma_{\mathrm{G}}(g)f\)（其中 \(g \in G\)）生成，是 \(\gamma_{G}\) 的有限维子表示。它等于其 \(\pi\)-同型分量的直和，对所有 \(\pi \in \widehat{G}\)（V, p. 457, 命题 1）。根据 V, p. 463, 推论 4，这意味着 \(E_{f}\) 包含于各空间 \(\varrho_{\mathrm{G}}(\pi)\) 的和中，从而 \(F_{c} \subset F_{b}\)。

用类似论证可得 \(F_{b} = F_{d}\)；由于 \(\varrho_{\mathrm{G}}(\pi)\) 是 \(\varrho_{G}\) 的子表示，同样可证明 \(F_{b} = F_{e}\)。

推论。--- 设 \(E \subset L^{2}(G)\) 是定义 \(\gamma_{G}\)（分别为 \(\delta_{G}\)、\(\varrho_{G}\)）的一个子表示的有限维向量子空间。则 E 包含于 \(\mathcal{C}(G)\)。

事实上，E 的每个元都是表示 \(\gamma_{G}\) 的 G-有限向量。

\subsection{分离拟完备空间中的表示}

命题 7. --- 设 \(\varrho\) 是 G 在分离、局部凸且拟完备空间 E 上的连续表示。E 的有限维子表示之和在 E 中稠密。

设 \(x \in \mathrm{E}\)，U 是 \(x\) 的一个开邻域。测度 \(\nu \in \mathcal{M}(\mathrm{G})\) 满足 \(\varrho(\nu)x \in \mathrm{U}\) 的所成集合，对于紧收敛拓扑在 \(\mathcal{M}(\mathrm{G})\) 中是开的（参见 V, p. 400, no. 2）。它包含 \(\varepsilon_e\)；因此它包含形如 \(\nu = f_1 \cdot \mu\) 的测度，其中 \(f_1 \in \mathcal{C}(\mathrm{G})\)（INT, VIII, p. 171, § 4, n° 7, 命题 19）；于是它还包含形如 \(f_2 \cdot \mu\) 的测度，其中 \(f_2\) 是 G-有限函数（V, p. 464, 命题 6 及 V, p. 462, 命题 4）。

\(\gamma_{G}\) 中由 \(f_{2}\) 生成的子表示 F 是有限维的。令 \(\widetilde{F}\) 为线性映射 \(f \mapsto \varrho(f)x\) 从 F 到 E 的像。空间 \(\widetilde{F}\) 是有限维的，且包含属于 U 的元 \(\varrho(f_{2})x\)。由于 \(\varrho(g)\varrho(f) = \varrho(\gamma_{\mathrm{G}}(g)f)\) 对所有 \(g \in G\) 和每个 \(f \in F\) 成立（V, p. 406, 公式 (1)），空间 \(\widetilde{F}\) 是 \(\varrho\) 的一个与 U 相交的子表示。命题得证。

推论 1. --- 设 π 是 G 在局部凸分离拟完备空间 E 上的连续不可约表示。空间 E 是有限维的。

推论 2. --- 设 \(\varrho\) 是 G 在局部凸分离拟完备空间 E 上的连续表示，\(\pi\) 是 G 的连续不可约表示。

\begin{enumerate}
\def\labelenumi{\alph{enumi})}
\item
  从 E 到 E 的 G-态射 \(p_{\pi} = \dim(\pi)\varrho(\overline{\chi}_{\pi})\) 是 E 的连续投影，其像是 \(\pi\)-同型分量的 \(\varrho\)；
\item
  若 \(\varrho\) 是酉表示，则投影 \(p_{\pi}\) 是 E 的正交投影，其像为 \(\mathrm{M}_{\pi}(\varrho)\)。
\end{enumerate}

证明 \(a\)。线性映射 \(p_{\pi} = \dim(\pi)\varrho(\overline{\chi}_{\pi})\) 定义良好，因为 E 拟完备且 G 紧（V, p. 401）。它属于 \(\mathrm{Hom}_{\mathrm{G}}(\varrho,\varrho)\)，因为 \(\pi\) 的特征标是 G 上的连续中心函数。

设 \(\varpi\) 是 E 的有限维子表示，F 是其空间。映射 \(p_{\pi}\) 通过对子空间的传递在 F 上诱导自同态 \(\dim(\pi)\varpi(\overline{\chi}_{\pi})\)。因此，当 \(\varpi\) 不同构于 \(\pi\) 时此自同态为零，否则是 F 的恒等映射（事实上，当 \(\varpi\) 不可约时由 V, p. 460, 推论 5 得到一般情形，而一般情形由此推出，因为由 V, p. 457, 命题 1，\(\varpi\) 是半单的）。

最后，由于 E 的有限维子表示之和在 E 中稠密（命题 7），且 \(p_{\pi}\) 连续，我们得到 \(p_{\pi}\) 是一个投影，其像是 \(\pi\)-同型分量的 \(\varrho\)。若 E 是希尔伯特空间，则由 V, p. 458, 推论 1，投影 \(p_{\pi}\) 是埃尔米特的，因而是正交投影（I, p. 133, 引理 3, (ii)）。

推论 3. --- 设 \(\varrho\) 是 G 在局部凸分离拟完备空间 E 上的连续表示。自同态

\[
x \mapsto \int_ {\mathrm{G}} \varrho (g) x d \mu (g)
\]

是 E 的一个投影，其像是 E 中不变向量的空间 \(E^{G}\)。特别地，若 E 有限维，则

\[
\mathrm{dim} (\mathrm{E} ^ {\mathrm{G}}) = \int_ {\mathrm{G}} \chi_ {\varrho} (g) d \mu (g).
\]

第一个断言是前一推论在 \(\pi\) 为 G 的一维平凡表示时的特殊情形。第二个断言由此得出，因为 \(E^{G}\) 的维数是到 \(E^{G}\) 上的正交投影的迹，即

\[
\mathrm{dim} (\mathrm{E} ^ {\mathrm{G}}) = \mathrm{Tr} \Bigl (\int_ {\mathrm{G}} \varrho (g) d \mu (g) \Bigr) = \int_ {\mathrm{G}} \chi_ {\varrho} (g) d \mu (g).
\]

特别地，当 E 有限维时，E 中存在向量 \(x \neq 0\)，使得 \(\varrho(g)x = x\) 对所有 \(g \in G\) 成立，当且仅当

\[
\int_ {\mathrm{G}} \chi_ {\varrho} (g) d \mu (g) \neq 0.
\]

推论 4. --- 设 \(\varrho\) 是 G 在希尔伯特空间 E 上的酉表示。空间 E 是各 \(\pi\)-同型分量 \(\mathrm{M}_{\pi}(\varrho)\)（其中 \(\pi \in \widehat{\mathrm{G}}\)）的希尔伯特直和。特别地，表示 \(\varrho\) 是半单的。

对应于 G 的不同构不可约表示的 \(\varrho\) 的同型分量彼此正交（V, p. 394, 命题 8）。由于 G 的每个有限维酉表示都是半单的（V, p. 457, 命题 1），各同型分量 \(M_{\pi}(\varrho)\) 对 \(\pi \in \widehat{G}\) 的和在 E 中稠密（命题 7）。推论由此得出。